
\documentclass[12pt,letterpaper]{article}
\usepackage{amssymb}
\usepackage{amscd}
\usepackage{amsmath}
\usepackage{amsfonts}
\usepackage{newtxtext,newtxmath}

\newtheorem{theorem}{Theorem}
\newtheorem{corollary}[theorem]{Corollary}
\newtheorem{conjecture}[theorem]{Conjecture}
\newtheorem{lemma}[theorem]{Lemma}
\newtheorem{proposition}[theorem]{Proposition}
\newtheorem{definition}[theorem]{Definition}
\newtheorem{example}[theorem]{Example}
\newtheorem{axiom}{Axiom}
\newtheorem{remark}{Remark}
\newtheorem{exercise}{Exercise}[section]

\def\ppp{{\mathbb{P}}}
\def\hhh{{\cal H}}
\def\ddd{{\cal D}}
\def\bbb{{\cal B}}
\def\sss{{\cal S}}
\def\zz{{\cal Z}}
\def\aaa{{\mathbb{A}}}
\def\fff{{\mathbb{F}}}
\def\qqq{\mathbb{Q}}
\def\rrr{\mathbb{R}}
\def\ccc{\mathbb{C}}
\def\zzz{\mathbb{Z}}
\def\nnn{\mathbb{N}}
\def\pf{{\bf proof}:\ }
\def\qed{$\Box$}
\def\eps{\epsilon}
\def\okk{{\cal O}_K}
\def\off{{\cal O}_F}
\def\oee{{\cal O}_E}
\def\ovg{\overline{g}}
\def\tarrow{\stackrel{\sim}{\rightarrow}}
\def\ovF{\overline{F}}
\def\ovx{\overline{x}}
\def\ovt{\overline{t}}
\def\bbg{{\cal B}(G)}
\def\bbh{{\cal B}(H)}
\def\bbH1{{\cal B}(H_1)}
\def\bbh2{{\cal B}(H_2)}
\def\sne{\subseteq}
\def\tr{{\rm tr}\,}

\begin{document}


\author{David Joyner}
\title{On twisted orbital integral identities\\ for $PGL(3)$
over a p-adic field\thanks{Appeared in Canadian Journal of Mathematics,
vol 42, 1990, pp. 1098-1130. Latexed 3-2003, 
correcting minor typos, with some minor changes to
grammar and the construction of the figures.
Research partially supported by an NSF fellowship, the Mathematics Department
of Princeton University and the Institute for Advanced Study.}}
\date{3-14-2003}
\maketitle

\begin{abstract}
The object of this paper is to prove certain $p$-adic 
orbital integral identities
needed in order to accomplish the symmetric square transfer via the 
twisted Arthur trace formula. Only \S \ref{sec:5} of this article 
contains original material, the rest of it
is due to R. Langlands. Very briefly,
we reduce the problem of proving certain
orbital integral identities for ``matching'' functions 
in the respective Hecke algebras to two counting problems on the buildings.
We give Langlands' solution of one of these problems in the case of the unit
elements of the respective Hecke algebras and \S \ref{sec:5}
provides the solution to the other one, again, in the unit element case. 
The main results assume $p\not= 2$.

\end{abstract}


\tableofcontents

\newpage

% === BEGIN ORIGINAL adjoint_lift1.tex ===

\section{Contents}
\label{sec:0}

Only \S \ref{sec:5} 
of this article contains original material, the rest of it
can be essentially be found in the unpublished notes \cite{La1}. 
Very briefly, \S \ref{sec:1} provides an 
introduction including the definition (due to
H. Jacquet and J. Shalika) of the ``norm map'', 
\S \ref{sec:2} applies the Satake transform
to proving some simple orbital integral identities in the ``split case'',
\S \ref{sec:3} reduces the problem of 
proving these identities in the ``non--split
case'' to two counting problems on the buildings, 
\S \ref{sec:4} recalls
R. Langlands' solution of one of these problems for the unit 
element of the Hecke algebra, and, finally, 
\S \ref{sec:5}
provides the solution to the other one (for the unit 
element of the Hecke algebra). 
In sections \ref{sec:4} and \ref{sec:5} we assume that $p\not= 2$.
\par
This paper was written independently of a recent paper \cite{F}
which also relied on Langlands' unpublished manuscipt \cite{La1}.
\vskip .1in

\noindent
Acknowledgement: The author sincerely thanks Professor R. Langlands for his
suggestions, many helpful conversations on the subject, and the use of
his notes \cite{La1}. Conversations with R. Kottwitz were also beneficial. 
Finally, I'd like to thank the referee
for some helpful comments. Most of this work was carried out at the
Institute for Advanced Study in 1986-7, supported by an NSF fellowship. 

\section{Introduction}
\label{sec:1}

\subsection{History}

Let $F$ denote a $p$--adic field and $\Pi (G(F))$ denote the set of
equivalence classes of admissible irreducible representations of $G(F)$,
for a 
reductive algebraic group $G$ defined over $F$. 
Typically, $F$ will be a completion of a
number field $K$ at a place $v$; let $\aaa = \aaa_K$ denote the adele ring 
of $K$.
 
One example of the local functoriality conjecture (\cite{Bor}, \cite{La2})
says that the 
$L$--homomorphism

\begin{equation}
\label{eqn:1.1}
r : GL(2,\ccc ) \rightarrow GL(2n+1,\ccc ),
%\tag 1.1
\end{equation}
given by $r := \text{Sym}^{\otimes 2n}\otimes (\text{det})^{-n}$ should 
yield a transfer (the ``symmetric $n^{th}$ power lifting'')

\begin{equation}
\label{eqn:1.2}
r_* : \Pi (GL(2,F)) \rightarrow \Pi (GL(2n+1,F)) .
%\tag 1.2
\end{equation}
Using $L$--parameters one can directly ``lift'' or transfer unramified
principal series from $GL(2,F)$ to $GL(2n+1,F)$, so there is some evidence
for the validity of this conjecture. There is of course a global analog
of this conjectural transfer, predicted by the global functoriality
principle. 
 
If $n=1$ then S. Gelbart and H. Jacquet \cite{GJ} proved the existence
of a global representation $\Pi \in \Pi (GL(3,\aaa ))$ associated
to a cuspidal $\pi \in \Pi (GL(2,\aaa ))$ using the theory of $L$--functions,
converse Hecke theory on $GL(3,\aaa )$ \cite{JPSS}, and an idea of
G. Shimura \cite{Sh}. Let me reformulate this in terms of functoriality: 
when $n = 1$, the $L$--map (\ref{eqn:1.1}) becomes

\begin{equation}
\label{eqn:1.3}
r\left(\smallmatrix a&b \\ c&d \endsmallmatrix\right)
= {1\over{ad-bc}}\left(\smallmatrix a^2&-ab&-b^2 \\
-2ac&ad+bc&2bd\\
-c^2&cd&d^2 \endsmallmatrix\right)
%\tag 1.3
\end{equation}
the adjoint representation of $sl(2,\ccc )$ with respect to the basis

\[
\underline{e}_1 := \left(\smallmatrix 0&0 \\ 1&0 \endsmallmatrix\right),
\ \underline{e}_2 := \left(\smallmatrix 1&0 \\ -1&0 \endsmallmatrix\right),
\ \underline{e}_3 := \left(\smallmatrix 0&1 \\ 0&0 \endsmallmatrix\right),
\]
and this factors through $PGL(2)$ giving us

\begin{equation}
\label{eqn:1.4}
\lambda_1 :  {^{L}SL(2,\ccc )} =PGL(2,\ccc )
\hookrightarrow SO(3,\ccc ) 
\hookrightarrow  SL(3,\ccc ) = {^{L}PGL(3,\ccc )} .
%\tag 1.4
\end{equation}
The functoriality conjecture then predicts a local and global transfer
from $SL(2)$ to $PGL(3)$. On the other hand, the natural embedding gives us
an $L$--map

\begin{equation}
\label{eqn:1.5}
\lambda_2 : {^{L}PGL(2,\ccc )} \hookrightarrow {^{L}PGL(3,\ccc )},
%\tag 1.5
\end{equation}
In this case the functoriality conjecture predicts a local and global
transfer from $PGL(2)$ to $PGL(3)$. The global transfer associated to 
(\ref{eqn:1.5}) 
is in fact the Gelbart--Jacquet lift. The local transfer associated to
(\ref{eqn:1.4}) will be elaborated upon in subsection 
\ref{sec:1.2} below.
 
Around 1976, H. Jacquet suggested that there should be a twisted trace
formula for $PGL(3)$, corresponding to the outer automorphism

\[
\sigma : g \longmapsto J{^{t}g^{-1}}J,\ \ \sigma^2 = 1,
\]
where $J$ is the idempotent matrix given by

\[
J:= \left(\smallmatrix \ &\ &1\\
\ &-1&\ \\
1&\ &\ \endsmallmatrix\right),
\]
and that this trace formula could be applied to proving the symmetric square 
transfer in a fashion analogous to the base change lift of 
Saito--Shintani--Langlands.  Moreover,
$PGL(2)$ and $SL(2)$ should be $\sigma$--endoscopic groups for $PGL(3)$ and the
symmetric square transfer should be a consequence of stabilizing this twisted
$PGL(3)$-trace formula. The notes \cite{La1}, completely a year or so later,
were motivated by Jacquet's suggestion.

\subsection{$L$--parameters and the transfer from $SL(2)$ to $PGL(3)$}
\label{sec:1.2}

Let $F$ be a $p$--adic field with ring of integers $\off$, let $G$
denote one of the groups $SL(n)$ or $PGL(n)$, let $K$ denote
a maximal compact subgroup of $G(F)$, and let $\hhh (G,K)$
denote the {\bf Hecke algebra} of compactly
supported locally constant functions on the double coset space
$K\setminus G(F)/K$. Let $K_0:= G(\off )$ denote the hyperspecial
maximal compact subgroup; when $K = K_0$, we write $\hhh (G) := 
\hhh (G,K_0)$. Let $C^\infty_c(G)$ denote the algebra of 
{\bf Schwartz--Bruhat functions} of locally constant functions of
compact support on $G(F)$. 
 
Corresponding to the $L$--map (\ref{eqn:1.4}), the Satake transform (defined 
in \S \ref{sec:2} below) gives us a homomorphism of Hecke algebras

\begin{equation}
\label{eqn:1.6}
\begin{split}
&\lambda_1^* : \hhh (PGL(3)) \rightarrow \hhh (SL(2)),\\
&\ \ \ \ \ \ \ \ \ \ \ \phi \longmapsto f :=\lambda_1^*(\phi),
\end{split}
%\tag 1.6
\end{equation}
determined by 
$\hat{f}(t) = \hat{\phi }(\lambda_1 (t)),\ t \in {^{L}SL(2,\ccc )}$.
When this identity holds, one says that the spherical function $\phi $
on $PGL(3,F)$ {\bf corresponds} with the spherical function $f$ on $SL(2,F)$.
Similarly, corresponding to the $L$--map (\ref{eqn:1.5})
there is the $L$--homomorphism of Hecke algebras

\begin{equation}
\label{eqn:1.7}
\begin{split}
&\lambda_2^* : \hhh (PGL(3)) \rightarrow \hhh (PGL(2)),\\
&\ \ \ \ \ \ \ \ \ \ \ \phi \longmapsto f := \lambda_2^*(\phi),
\end{split}
%\tag 1.7
\end{equation}
determined by $\hat{f}(t) = \hat{\phi }(\lambda_2 (t)),\ t \in 
{^{L}PGL(2,\ccc )}$. Using this one can define a ``correspondence'', as above,
between spherical functions on $PGL(3,F)$ and spherical functions on
$PGL(2,F)$. 
 
According to the local Langlands correspondence for $GL(3)$, to each
unramified admissible homomorphism
$\psi :W_F \rightarrow {^{L}SL(3,\ccc )}$ there is an $L$--packet of
admissible irreducible 
unramified representations $\Pi (\theta )  \in
\Pi (SL(2,F))_L$ of SL(2,F). Here $W_F$ denotes the 
absolute local Weil group of $F$ and the subscript $L$ on the $\Pi$ signifies
that we are identifying $L$--indistinguishable representations. 
Each homomorphism
$\theta : W_F \rightarrow {^{L}SL(2,\ccc )}$ yields via (\ref{eqn:1.4})
a homomorphism $\theta^* : W_F \rightarrow {^{L}PGL(2,\ccc )}$, and if 
$\theta$ is
admissible then so is $\theta^*$. The local functoriality conjecture
predicts that the $L$--map (\ref{eqn:1.4}) 
should yield a map of tempered $L$--packets

\begin{equation}
\label{eqn:1.8}
\begin{split}
&\Lambda_1 : \Pi^{\text{temp}}(SL(2,F))_L \rightarrow 
\Pi^{\text{temp}}(PGL(3,F)),\\
&\quad\quad \ \ \ \ \ \ \ \ \ \ \Pi (\theta ) \longmapsto \pi_{\theta^*}.
\end{split} 
%\tag 1.8
\end{equation}
We say that $\pi \in \Pi (SL(2,F))$ {\bf transfers} to 
$\Pi \in \Pi (PGL(3,F))$ if $\pi \in \Pi (\theta )$ and 
$\Pi = \pi_{\theta^*}$, for some admissible $\theta$ as above. 
The $L$--packets of $SL(2)$ have been described by J.-P. Labesse
and R. Langlands \cite{LL}.

\subsection{The Jacquet--Shalika norm maps}
\label{sec:1.3}

The basic ideas of this subsection are, I believe, due to H. Jacquet
and J. Shalika and may be found in \cite{La1}.
 
Let $G := PGL(3),\ H_1 := SL(2),\ H_2 := PGL(2)$.
Associated to the two $\sigma$-endoscopic groups $H_i$ are ``norm'' maps

\begin{equation}
\label{eqn:1.9}
N_i : \{\text{stable $\sigma$-conjugacy classes in G}\} \rightarrow
\{\text{stable conjugacy classes in $H_i$}\},
%\tag 1.9
\end{equation}
which (hopefully) allow one to relate stable class functions on $G$ to stable
class functions on $H_i$. The explicit construction of these norm maps
is the object of this subsection.

\begin{lemma}
\label{lemma:1.10} 
If $[A] \in PGL(2n+1,F)$ is represented by
$A \in GL(2n+1,F)$ then $[A]\sigma ([A])$ is represented by 
$A\cdot {^{t}A^{-1}}$
and $A\cdot {^{t}A^{-1}}$ has at least one eigenvalue equal to one.
\end{lemma}
 
This is an immediate consequence of the fact that $A^{-1} - {^{t}A^{-1}}$
is singular (since $2n+1$ is odd), hence has $0$ as an eigenvalue.
 
Let $V:=F^{2n+1}$ and identify, by fixing a basis of $V$, $GL(2n+1,F)$ with 
$\text{Aut}_F(V)$. For the $[A]$ and $A$ in the lemma, let 
$V_A \subset V$ be a 1-dimensional subspace on which $A\cdot {^{t}A^{-1}}$
acts as the identity. Let $W_A$ be the orthogonal complement of
${^{t}A^{-1}}V_A$ with respect to the inner product 
$(v_1,v_2) := {^{t}v_1\cdot v_2}$.
Observe that $A\cdot {^{t}A^{-1}}$ acts on $V_A$ trivially, 
therefore on $V/V_A$, and
$A{^{t}A^{-1}}$ acts on $W_A$. 
 
Let me now restrict to the case $n = 1$.

\begin{lemma}
\label{lemma:1.11} 
Suppose that $A\cdot {^{t}A^{-1}}$ induces the linear transformation
$B_1 \in \text{Aut}_F(W_A)$ and $B_2 \in \text{Aut}_F(V/V_A)$. Then

\begin{itemize}

\item[(a)] 
$\det B_1 = \det B_2 = 1$,

\item[(b)] 
there is an isomorphism $\theta : W_A \rightarrow V/V_A$ such that

\[
\begin{CD}
W_A @>B_1>>W_A\\
@V\theta VV @AA\theta A\\
V/V_A @>B_2>>V/V_A
\end{CD}
\]
commutes,

\item[(c)] 
as elements of $SL(2,F)$, $B_1$ and $B_2$ are stably conjugate,

\item[(d)] 
the stable conjugacy class of $B_i$ is independent of the choice 
of $V_A$.
\end{itemize}
\end{lemma}
 
For the proof, I refer to \cite{La1}.

\begin{definition}
\label{def:1.12} 
Choose a basis of $V/V_A$ or of $W_A$ and represent
the transformation induced by $A\cdot {^{t}A^{-1}}$ by a 
$2\times 2$ matrix. The stable
$SL(2)$-conjugacy class of this matrix is 
well--defined and depends only on the 
stable $\sigma$--conjugacy class of $[A]$. 
This stable conjugacy class in
$SL(2,F)$ is called the {\bf $H_1$-norm} of $[A]$, written as
$N_1([A]) \subseteq H_1(F)$. For the $H_2$--norm, first identify
$H_2$ with SO(3) via the adjoint representation (\ref{eqn:1.3}), 
where $SO(3)$ denotes the connected component of the orthogonal group for the 
skew--diagonal matrix $J$ above. If the eigenvalues of
$A\cdot {^{t}A^{-1}}$ are $\alpha,\ 1,\ \alpha^{-1}$ and 
$\alpha \not= \pm 1$, then
define the {\bf $H_2$--norm} of $[A]$, written $N_2([A])$, to be the
stable $H_2$--conjugacy class which contains a matrix with
$\alpha,\ 1,\ \alpha^{-1}$ as eigenvalues.
\end{definition}
 
As a matter of notation, let

\[
\begin{split}
&Cl^{st}_{\sigma}(G) :=\{\text{stable $\sigma$-conjugacy classes in G}\},\\
&Cl^{st}(H_i) := \{\text{stable conjugacy classes in $H_i$}\}.
\end{split}
\]

\begin{proposition} (Jacquet--Shalika)
\label{prop:1.13} 
The norm map

\[
N_1 : Cl^{st}_{\sigma}(G) \rightarrow Cl^{st}(H_1)
\]
is a bijection.
\end{proposition}

\begin{remark} 

(1) This is proven in \cite{La1} by explicit
matrix manipulations.
\newline

(2) For $H_2$ stable conjugacy and ordinary conjugacy
are essentially the same, so the statement analogous to the 
proposition is false for $N_2$. 
\newline

(3) In general, one expects to be able to define norms from PGL(r) to
any of its $\sigma$--endoscopic groups. One expects that the norm map
associated to the ``largest'' $\sigma$--endoscopic group should also yield a
bijection. For r = 2n+1 and r = 2n, one would like bijections

\[
\begin{split}
&N: Cl^{st}_{\sigma}(PGL(2n+1)) \rightarrow Cl^{st}(Sp(2n)),\\
&N: Cl^{st}_{\sigma}(PGL(2n)) \rightarrow Cl^{st}(PSp(2n)).
\end{split}
\]
\end{remark}
 
The following corollary of Langlands' proof of (\ref{prop:1.13}) 
is sometimes useful.

\begin{corollary}
\label{cor:1.14} 
If $N_1([A])$ does not contain a unipotent, then 
there is a $[A']$ in the same stable $\sigma$--conjugacy class 
as $[A]$ and 
represented by $A' \in GL(3,F)$, for which

\begin{itemize}

\item[(a)] 
$V = V_{A'} \oplus W_{A'}$,

\item[(b)] 
$W_{A'}^{\bot} = V_{A'}$ with repect to the inner product
$(v_1,v_2):={^{t}v_2}\cdot v_1$,

\item[(c)] 
$V_{A'}$ and $W_{A'}$ are both invariant under $A$ and 
${^{t}A^{-1}}$ (i.e., $A$ can be put in $(2,1)$-block form).
\end{itemize}
\end{corollary}


\subsection{The fundamental lemmas}
\label{sec:1.4}

Recall $G := PGL(3),\ H_1 := SL(2),\ H_2 := PGL(2)$. Let us denote the 
{\bf twisted centralizer} by

\begin{equation}
\label{eqn:1.15}
G(g\sigma ,F) := \{g \in G(F)\ |\ x^{-1}g\sigma (x) = g\},
%\tag 1.15
\end{equation}
and the (ordinary) {\bf centralizer} by

\begin{equation}
\label{eqn:1.16}
H_i(h,F) := \{h \in H_i(F) \ |\ x^{-1}hx = h\},\ \ \ i=1.2.
%\tag 1.16
\end{equation}
Let $\omega_{G(g\sigma )},\ \omega_G$ be fixed nonzero forms of maximal
degree on $G(g\sigma ,F),\ G(F)$, respectively, and let

\[
d\ovx := {d\omega_G(x)\over{d\omega_{G(g\sigma )}(x)}}
\]
denote the quotient measure. For $\phi \in C^\infty_c(G)$, define the 
{\bf twisted orbital integral} by

\begin{equation}
\label{eqn:1.17}
\Phi (g,\phi ,\omega_G,\omega_{G(g\sigma )}) = \Phi (g,\phi )
:= \int_{G(g\sigma ,F)\setminus G(F)}
\phi (g^{-1}g\sigma (g))\ d\ovg,
%\tag 1.17
\end{equation}
provided it converges.
Similarly, let $\omega_{H_i(h)},\ \omega_{H_i}$ be fixed nonzero 
forms of maximal
degree on $H_i(h\sigma ,F),\ H_i(F)$, respectively, and let

\[
d\ovx := {d\omega_{H_i}(x)\over{d\omega_{H_i(h)}(x)}}
\]
denote the quotient measure. In fact, we don't take any $\omega_{H_i(h)}$,
but only one obtained by pulling back $\omega_{G(g\sigma)}$ via an \'etale
surjective homomorphism over $\overline{F}$

\[
\psi_i : H_i(h)^0 \rightarrow G(g\sigma ),
\]
namely, it must satisfy

\[
\omega_{H_i(h)}=|\text{ker }\psi_i|^{-1}\psi^*_i\omega_{G(g\sigma )},
\]
for all $h \in N_i(g)$. Here $N_i$ denotes the norm map defined above. 
Note that if $h \in H_i(F)$ is regular semi--simple then we may assume that
$\psi_i$ is an isomorphism. For $f \in C^\infty_c(H_i)$ and $h \in H_i(F)$
regular, define the {\bf orbital integral} by

\begin{equation}
\label{eqn:1.18}
\Phi (h,f,\omega_{H_i},\omega_{H_i(h)}) = 
\Phi (h,f):= \int_{H_i(h,F)\setminus H_i(F)}
f(x^{-1}hx)\ d\ovx .
%\tag 1.18
\end{equation} 
This defines the orbital integrals which will
occur in the fundamental lemmas below. Note that, whereas 
$\Phi (h,f)$ is a class function on $H_i(F)^{reg}$, $\Phi (g,\phi )$
may be regarded as a class function on the nonconnected reductive group

\[
\overline{G} := G \rtimes \{1,\sigma \},\ \ {\rm (semi-direct\ product)},
\]
because of the identity

\[
(x\times \sigma )^{-1}(g\times \sigma )(x\times \sigma )
=\sigma (x^{-1}g\sigma (x))\times \sigma .
\]
 
To analyse these integrals further we of course should know something
about the conjugacy classes on $H_i$ and on $\overline{G}$. Actually,
from now on I shall pretty much ignore $\overline{G}$ since it
seems easier to regard conjugacy classes on
$\overline{G}$ as $\sigma$--conjugacy classes on $G$. 
 
Let me leave aside the notion of $\sigma$-conjugacy for the moment
and first discuss the (untwisted) notions of conjugacy and
stable conjugacy. Two elements
$h_1,\ h_2 \in SL(2,F) = H_1(F)$ can be {\bf stably conjugate}
(meaning that there is an $x \in H_1(\ovF )$ such that
$h_1 = x^{-1}h_2x$) but not conjugate (over $F$). 
This is not true for $PGL(n)$. 
In other words, if two elements of $PGL(n,F)$ are
$PGL(n,\ovF )$-conjugate then they are necessarily
$PGL(n,F)$-conjugate (see, for example, \cite{L}, ch. 15, \S 3, p. 543).
Equivalently, for $PGL(n)$, the notions of stable 
conjugacy and (ordinary) conjugacy coincide. 
For groups other than
$PGL(n)$, however, these notions do not generally coincide. 
As a measure of the 
difference between these two notions, one may consider the set
of conjugacy classes of an element within a given stable conjugacy class
of that element. 
For $h \in H_i(F)$, the set of $H_i(F)$-conjugacy
classes of $h_i$ within the stable conjugacy class of $h$ is parameterized 
by 

\begin{equation}
\label{eqn:1.19}
\ddd (h,H_i) := H_{i,h}(\ovF )\setminus St(h,H_i)/H_i(F),
%\tag 1.19
\end{equation}
where

\[
St(h,H_i) := \{x \in H_i(\ovF )\ |\ x^{-1}hx \in H_i(F)\}
\]
denotes those elements over $\ovF$ which leave invariant the
stable conjugacy class of $h \in H_i(F)$ under conjugation. 
If $i = 2$ then
$\ddd (h,H_2)$ is a singleton since in that case the notions of 
stable conjugacy and ordinary conjugacy coincide. 
 
Now let me turn to the twisted analogs of these notions. Although
$PGL(n)$ has ``no $L$--indistinguishabilty'' it \underbar{does} have
``twisted 
\newline
$L$--indistinguishability''. In other words, it is
possible for $g_1,\ g_2 \in G(F)$ to be stably $\sigma$-conjugate
(meaning that there is some $g\in G(\ovF )$ such that
$g_1 = g^{-1}g_2\sigma (g)$) but not $\sigma$-conjugate over $F$.
For $g \in G(F)$, the set of $\sigma$-conjugacy classes of $g$ within the
stable $\sigma$-conjugacy class of $g$ is parameterized by

\begin{equation}
\label{eqn:1.20}
\ddd_{\sigma}(g,G) := G_{g\sigma}(F)\setminus St_{\sigma}(g)/H_i(F),
%\tag 1.20
\end{equation}
where

\[
St_{\sigma}(g,G) := \{x \in G(\ovF )\ |\ x^{-1}gx \in G(F)\}
\]
denotes those elements over $\ovF$ which leave invariant the 
stable $\sigma$-conjugacy class of $g$.

\begin{lemma}
\label{lemma:1.21} 
There are isomorphisms (as pointed sets)

\[
\begin{split}
&\ddd (h, H_i) \cong H^1(F,H_i(h)),\\
&\ddd_{\sigma}(g,G) \cong H^1 (F,G(g\sigma)).
\end{split}
\]
\end{lemma}
 
For the proof of this see, for example, \cite{La2}. We will use this to
compute the cardinality of these sets later.
 
For fixed $g \in G(F)$ and arbitrary $g' \in \ddd_{\sigma}(g,G)$, 
all the twisted centralizers $G(g'\sigma)$ are (well--defined and)
isomorphic over $\ovF$. 
When considering the twisted orbital integrals
$\Phi (g',\omega_{G(g'\sigma )},\omega_G),\ g' \in \ddd_{\sigma}(g,G)$, 
we will always assume that the measures $\omega_{G(g'\sigma )}$ are determined
from the given $\omega_{G(g\sigma )}$ by pulling it back via some isomorphism
$G(g'\sigma) \tarrow G(g\sigma)$. 
A similar remark pertains to the 
orbital integrals attached to the $h' \in \ddd (h,H_i)$.
 
If $N_1(g)$ is regular, $g \in G(F)$ then define the {\bf stable twisted 
orbital integral} of $\phi \in C^\infty_c(G)$ by

\begin{equation}
\label{eqn:1.22}
\Phi^{T,1}(g,\phi ):= \sum_{g' \in \ddd_{\sigma}(g,G)} \Phi (g',\phi ),
%\tag 1.22
\end{equation}
where $T$ denotes $G(g\sigma )$.
Similarly, if $h \in H_i(F)$ is regular then define the
{\bf stable orbital integral} of $f \in C^\infty_c(H_i)$ by

\begin{equation}
\label{eqn:1.23}
\Phi^{T_i,1}(h,f):= \sum_{h' \in \ddd(h,H_i)} \Phi (h',f),
%\tag 1.23
\end{equation}
whete $T_i$ denotes the centralizer $H_i(h)$.

\begin{theorem} (Langlands,
{\bf Fundamental Lemma for $H_1$})
\label{thrm:fundlemmaH1}
If
$\phi \in \hhh (G)$ and $f_1 = \lambda_1^*(\phi ) \in
\hhh (H_1)$ then, for all non--trivial $h \in N_1(g)$, we have

\begin{equation}
\label{eqn:1.24}
\Phi^{T,1}(g,\phi ) = \Phi^{T_1,1}(h,f_1).
%\tag 1.24
\end{equation}
\end{theorem}
 
We will prove this in \S \ref{sec:4} 
in the case where $\phi$ is the unit element
of $\hhh (G)$ and $p\not= 2$.
 
The analogous statement for $H_2$ requires the introduction
of ``unstable'' orbital integrals. These are the same as stable orbital
integrals except that the sum is ``twisted'' by certain roots of unity
($\pm 1$'s in our case). To define these, we use the property that
the elements $g \in G(F)$ with $N_1(g) = 1$ have the property
that $|\ddd_{\sigma}(g,G)|=2$, and we can define a bijection of sets

\begin{equation}
\label{eqn:1.25}
\kappa : \ddd_{\sigma}(g,G) \rightarrow \{\pm 1\},
%\tag 1.25
\end{equation}
by

\[
\kappa (g) := 
\begin{cases} 
+1,&\ \ \ {\rm if}\ G(g\sigma)\ {\rm is\ split\ over}\ F,\\
-1,&\ \ \ {\rm if}\ G(g\sigma)\ {\rm is\ non--split\ over}\ F .
\end{cases}.
\]
If $N_1(g) \not= 1$ is regular then replace $g$ by $g_{s} :=
(g+{^{t}g})/2$ to get a regular element with $N_1(g_{s})=1$. In this
case, define $\kappa (g) = \kappa (g_{s})$. This depends only
on the $\sigma$-conjugacy class of $g$. This sign is closely related
to the Hilbert norm residue symbol of $E/F$, where $E$ is the quadratic
extension obtained from adjoining the eigenvalues of $h \in N_1(g)$ to
$F$, see \S \ref{sec:4} below. 
The {\bf unstable or $\kappa$-orbital integral}
of $\phi \in C^\infty_c(G)$ is

\begin{equation}
\label{eqn:1.26}
\Phi^{T,\kappa}(g,\phi ):= \sum_{g' \in \ddd_{\sigma}(g,G)} 
\kappa (g')\Phi (g',\phi ),
%\tag 1.26
\end{equation}
where $T:=G(g\sigma )$.

\begin{conjecture}
({\bf Fundamental Lemma for $H_2$})
Assume $N_2(g) \not= 
\pm 1$. If $\phi \in \hhh (G)$ and $f_2:=\lambda_2^*(\phi ) \in
\hhh (H_2)$ then, for any $h \in N_2(g)$, we have

\begin{equation}
\label{eqn:1.27}
\Phi^{T,\kappa}(g,\phi ) = \tau (g)\Phi(h,f_2),
\end{equation}
where $\tau$ is the transfer factor
\[
\tau (h) := \pm |(1+\beta_{1})(1+\beta_{2})|^{1/2},
\]
denoting the eigenvalues of $h$ by $\beta_{1},\ \beta_{2}$. The $\pm$ sign
will be explained in (\ref{eqn:3.18}) below.
\end{conjecture}
 
This main result of this paper is, as mentioned earlier, the proof that
this holds when the residual characteristic of $F$ is greater than $2$ and
$\phi \in \hhh (G)$ is the unit element. 






% === END ORIGINAL adjoint_lift1.tex ===
% === BEGIN ORIGINAL adjoint_lift2.tex ===

\section{The identities in the split case}
\label{sec:2}

\subsection{Summary}

The object is to briefly sketch, following \cite{La1}, a proof of the stable
and unstable fundamental lemmas

\begin{equation}
\label{eqn:2.1}
\Phi^{T,1}(g,\phi ) = \Phi^{T_1,1}(h,f),
%\tag 2.1
\end{equation}
for $h\in N_1(g)\cap A_1(F)$ regular semi--simple, and

\begin{equation}
\label{eqn:2.2}
\Phi^{T,\kappa}(g,\phi ) = |(1+\beta_1)(1+\beta_2)|^{-1/2}\Phi(h',f),
%\tag 2.2
\end{equation}
for $h'\in N_2(g)\cap A_2(F)$ regular, semi--simple. Here $A_i(F)$ denotes
the maximal split torus in $H_i(F)$; let $A(F)$ denote the maximal split torus
of G(F). The norms have been defined in such a way that, if $g\in A(F)$
then $N_i(g)\cap A_i(F) \not= \emptyset$. We use Satake transforms,
a method which has the advantage of yielding, for the split groups
we consider, an explicit expression for the absolute value of the
transfer factor (the precise transfer factor is very difficult to predict in
general, see the recent work of R. Langlands and D. Shelstad \cite{LS}).

\subsection{2.2. Background on twisted integration formulas}
\label{sec:2.2}

All Haar measures in this subsection are normalized as in \cite{C}.
\par
In this section, let $G' =SL(n)$ or $PGL(n)$ 
and consider the outer automorphism
$\sigma (g) = J'\cdot {^{t}g}^{-1}\cdot J'$ where $J'$ denotes the $n\times n$
skew--diagonal matrix with alternating $\pm 1's$ on the skew--diagonal:

\[
J':=
\left(
\begin{array}{cccccc}
 \ &\ &\ &\ &\ &1\\
\ &\ &\ &\ &-1&\ \\
\ &\ &\ &.&\ &\ \\
\ &\ &.&\ &\ &\ \\
\ &.&\ &\ &\ &\ \\
(-1)^{n-1}&\ &\ &\ &\ &\
\end{array}
\right) .
\]
Let $Ad$ denote the action of $G'(F)$ on Lie$(G')$:
$Ad(g)X :=gXg^{-1}$. 
The Jacobian of the change--of--variables
$n\longmapsto nmn^{-1}m^{-1} =\exp[(1-Ad(m))\log n]$,
where $\log n\in Lie\  N'$ is given by $\exp \log n :=n$, is
of course $\Delta (m) := |{\rm det} (Ad(m)-1)_{Lie\ N'}|_F $.
Similarly, the Jacobian of $n\longmapsto nm\sigma (n)^{-1}m^{-1}
=\exp [(1-Ad(m)\sigma )\log n]$ is $\Delta_{\sigma}(m) := |{\rm det}
(Ad(m)\sigma - 1)_{Lie\ N'}|_F$.
 
To be concrete, if $\alpha_1,\dots,\alpha_n$ are the simple roots of $G'$ then
$\sigma$ acts on their associated root groups via
$\sigma (\exp (tX_{\alpha_i})) = J(\exp (tX_{\alpha_i}))^{-1}J
= (-1)^{n-1}\exp (-tJ{^{t}X}_{\alpha_i}J)$,
or equivalently, via the action of $\sigma \in Out(G')\cong
{\rm Aut}(\ddd )$ on the Dynkin diagram $\ddd $ of $G'$. 
For $PGL(3)$, $\sigma$ exchanges the two simple roots; on the level of the
Lie algebra, it exchanges the root vectors $X_{\alpha_1},
\ X_{\alpha_2}$ and changes the sign of $X_{\alpha_1+\alpha_2}$
in Lie$\,N'$ since $X_{\alpha_1+\alpha_2}=[X_{\alpha_1},X_{\alpha_2}]$.
Therefore the analog to

\begin{equation}
\label{eqn:2.3}
\int_{N'} \phi (nmn^{-1}m^{-1})\ dn
= \Delta (m)^{-1}\int_{N'} \phi (n)\ dn 
%\tag 2.3
\end{equation}
becomes

\begin{equation}
\label{eqn:2.4}
\int_{N'} \phi (nm\sigma (n)^{-1}m^{-1})\ dn
= \Delta_{\sigma}(m)^{-1}\int_{N'} \phi (n)\ dn 
%\tag 2.4
\end{equation}
 
Let $a\in A'(F)$
belong to the maximal split torus of $G'(F)$ and let $\phi$ be a spherical
function on $G'(F)$. As a straightforward application of 
(\ref{eqn:2.3})-(\ref{eqn:2.4}), we
have

\begin{equation}
\label{eqn:2.5}
\begin{split}
\int_{G'(g\sigma ,F)\setminus G'(F)}
&\phi (x^{-1}g\sigma (x))\ d\ovx\\
&=\int_{G'(g\sigma ,F)\setminus A'(F)}\Delta_{\sigma}(m)^{-1}\Delta (m)
\int_{A'(F)\setminus G'(F)} \phi (x^{-1}mx)\ d\ovx d\ovt ,
\end{split}
%\tag 2.5
\end{equation}
where $m:=t^{-1}a\sigma (t)$. Notice that the outer integral on the
right is a finite sum. This fact will be used in a later subsection.
 
The remainder of this subsection is devoted to the calculation of the
twisted Jacobian $\Delta_{\sigma}(m)$, when $G' = G = PGL(3)$. Write
Lie$\,N = F\cdot X_{\alpha_1}\oplus F\cdot X_{\alpha_2}
\oplus F\cdot X_{\alpha_1+\alpha_2}$ and identify its elements as
column vectors. In terms of this basis, $\sigma$ is represented by
the matrix

\[
\left( 
\begin{array}{ccc} 
\ &-1&\ \\
-1&\ &\ \\
\ &\ &-1
\end{array} 
\right)
\]
and $Ad(a)$ is represented by the matrix

\[
\left( 
\begin{array}{ccc} 
\alpha_1(a)&\ &\ \\
\ &\alpha_2(a)&\ \\
\ &\ &(\alpha_1+\alpha_2)(a)
\end{array} 
\right),
\]
where of course $\alpha_1(a) = a_1/a_2,\ \alpha_2(a) = a_2/a_3,
\ (\alpha_1 +\alpha_2)(a) = a_1/a_3$. Combining this together, we find
that ${\rm det} (Ad(a)\sigma -1)_{Lie\, N}= (a_1/a_3)^2-1$.
In particular, with $a = m = t^{-1}g\sigma (t) \left( \smallmatrix
\gamma_1t_1^{-1}t_3^{-1}&\ &\ \\
\ &\gamma_2t_2^{-2}&\ \\
\ &\ &\gamma_3t_1^{-1}t_3^{-1}\endsmallmatrix\right)$,
where $\gamma_i$ denote the eigenvalues of a representative of $g\in A(F)$,
we find that

\begin{equation}
\label{eqn:2.6}
\Delta_{\sigma}(m) = |1-(\gamma_1/\gamma_3)^2|_F.
%\tag 2.6
\end{equation}

\subsection{The Satake transform on $PGL(3)$, $SL(2)$, $PGL(2)$}
\label{sec:2.3}

There are, as we mentioned already in \S \ref{sec:1}, $L$--maps

\[
\begin{split}
&\lambda_1:{^{L}H_1}(\ccc ) \rightarrow {^{L}G}(\ccc ),
\ \ \ {^{L}H_1} = SO(3)\\
&\lambda_2:{^{L}H_2}(\ccc ) \rightarrow {^{L}G}(\ccc ),
\ \ \ {^{L}H_2} = SL(2),
\end{split}
\]
given on the maximal split torus by

\[
\begin{split}
&\lambda_1 : 
\left( 
\begin{array}{cc} 
t_1&\ \\ 
\ &t_2
\end{array} 
\right) 
\longmapsto
\left( 
\begin{array}{ccc} 
t_1t_2^{-1}&\ &\ \\
\ &1&\ \\
\ &\ &t_2t_1^{-1}
\end{array} 
\right) ,\\
&\lambda_2 : 
\left( 
\begin{array}{ccc} 
t&\ \\ \ &t^{-1}
\end{array} 
\right) \longmapsto
\left( 
\begin{array}{ccc} 
t&\ &\ \\
\ &t^{-1}&\ \\
\ &\ &1
\end{array} 
\right) .
\end{split}
\]
Define the {\bf Satake transform} by

\begin{equation}
\label{eqn:2.7}
\begin{split}
&\hat{\phi}(t) := \sum_{\Lambda \in X^*({^{L}A})} \alpha (\Lambda )\Lambda (t),
\ \ t\in {^{L}A}(\ccc ),\\
&\hat{f_1}(t) := \sum_{\mu \in X^*({^{L}A_1})} \beta_1(\mu )\mu (t),
\ \ t\in {^{L}A_1}(\ccc ),\\
&\hat{f_2}(t) := \sum_{\nu \in X^*({^{L}A_2})} \beta_2(\nu )\nu (t),
\ \ t\in {^{L}A_2}(\ccc ).
\end{split}
%\tag 2.7
\end{equation}
Here

\begin{equation}
\label{eqn:2.8}
\alpha (\Lambda ):=\Delta_G (g)\int_{A(F)\setminus G(F)}
\phi (x^{-1}gx)\ d\ovx ,
%\tag 2.8
\end{equation}
where $\Lambda$ is that element $\Lambda =\Lambda (g)
=\Lambda (k_1,k_2,k_3) \in X^*({^{L}A})
\cong \zzz^3$ is that element corresponding to $(k_1,k_2,k_3)\in \zzz^3$
for $k_i$ given by

\[
g = \left( \begin{array}{ccc} u_1\pi^{k_1}&\ &\ \\
\ &u_2\pi^{k_2}&\ \\
\ &\ &u_3\pi^{k_3}\end{array} \right) ,\ \ u_i\in \off^{\times}.
\]
We also define $\Delta_G(g)
:=\prod_{{\rm all\  roots }\alpha }|1-\alpha (g)|^{1/2}
=|1-\gamma_1/\gamma_3||1-\gamma_2/\gamma_3||1-\gamma_1/\gamma_3|
|\gamma_3/\gamma_1|.$ Similarly, for $H_1(F)$ we have

\begin{equation}
\label{eqn:2.9}
\eta_1 (\mu ):=\Delta_{H_1} (h)\int_{A_1(F)\setminus H_1(F)}
f_1(x^{-1}hx)\ d\ovx ,
%\tag 2.9
\end{equation}
where $\mu =\mu (h)= \mu (m) \in X^*({^{L}H_1})\cong \zzz$ corresponds to

\[
h = \left( \begin{array}{cc} \beta_1&\ \\ \ &\beta_2
\end{array}
\right)
= \left( 
\begin{array}{cc}
u'_1\pi^{m}&\ \\ 
\ &u'_2\pi^{-m}
\end{array}
\right),
\]
where $u'_i\in \off^{\times}$.
Here we define $\Delta_{H_1}(h)
:=\prod_{{\rm all\  roots }\alpha }|1-\alpha (h)|^{1/2}
=|1-\beta_1/\beta_2||\beta_2/\beta_1|^{1/2}.$
For $H_2(F)$ we define $\Delta_{H_2}(h')
:=\prod_{{\rm all\  roots }\alpha }|1-\alpha (h')|^{1/2}
=|1-\beta '_1||\beta '_1|^{-1/2}$, and

\begin{equation}
\label{eqn:2.10}
\eta_2 (\nu ):=\Delta_{H_2} (h')\int_{A_2(F)\setminus H_2(F)}
f_2(x^{-1}h'x)\ d\ovx ,
%\tag 2.10
\end{equation}
where $\nu = \nu (n_1,n_2) \in X^*({^{L}H_2})$ corresponds to

\[
h' = \left( 
\begin{array}{cc}
\lambda\beta '_1&\ \\ \ &\lambda
\end{array}
\right)
= \left( 
\begin{array}{cc}
u''_1\pi^{n_1}&\ \\ \ &u''_2\pi^{n_2}
\end{array}
\right) ,
\]
where
$u''_i \in \off^{\times}$. Here the pair
$(n_1,n_2) \in\zzz^2$ is not well--defined but $n_1 -n_2$ is.
 
Note that, via the induced map
$\lambda_1^* : X^*({^{L}A})\rightarrow X^*({^{L}A_1}),$
two $\Lambda (k_1,k_2,k_3)$, 
$\Lambda (k'_1,k'_2,k'_3)$ belonging to $X^*({^{L}A})$
have the same image if and only if $k_1-k_3 = k'_1-k'_3$. Similarly,
via the induced map
$\lambda_2^* : X^*({^{L}A})\rightarrow X^*({^{L}A_2}),$
two $\Lambda (k_1,k_2,k_3)$, 
$\Lambda (k'_1,k'_2,k'_3)$ in $X^*({^{L}A})$
have the same image if and only if $k_1+k_2 = k'_1+k'_2=2k_3=2k'_3$.

\begin{lemma}
\label{lemma:2.11}

\begin{itemize}
\item[(a)] 
For $f_1$ and $\phi$ spherical, the following are equivalent:

\begin{itemize}
\item[(i)] 
$f_1=\lambda_1^*\phi$,

\item[(ii)] 
$\hat{f_1}(t) = \hat{\phi}(\lambda_1(t)),\ \ t\in {^{L}H_1}$,

\item[(iii)]

\[
\sum_{\substack{\Lambda = \Lambda (k_1,k_2,k_3)\\k_1-k_3=m}} 
\alpha (\Lambda )
=\eta_1 (\mu ),\ \ \forall \mu = \mu (m).
\]
\end{itemize}
\item[(b)] 
For $f_2$ and $\phi$ spherical, the following are equivalent:

\begin{itemize}
\item[(i)] 
$f_2=\lambda_2^*\phi$,

\item[(ii)]
$\hat{f_2}(t) = \hat{\phi}(\lambda_2(t)),\ \ t\in {^{L}H_2}$,

\item[(iii)]

\[
\sum_{\substack{\Lambda = \Lambda (k_1,k_2,k_3)\\k_1=n_1,\ k_2=n_2}}
\alpha (\Lambda )
=\eta_2 (\nu ),\ \ \forall \nu = \nu (n_1,n_2).
\]
\end{itemize}
\end{itemize}
\end{lemma}
 
The proofs are omitted.

\subsection{The fundamental identities in the split case}
\label{sec:2.4}

We may always
represent $t\in A(F)$ by a diagonal matrix in $GL(3,F)$ whose middle
entry $t_2$ is $1$. This gives

\[
t^{-1}\sigma (t)
=
\left( 
\begin{array}{ccc} 
t_1^{-1}t_3^{-1}&\ &\ \\
\ &1&\ \\
\ &\ &t_1^{-1}t_3^{-1}
\end{array} 
\right) .
\]
Because of this, the right side of the expression in (\ref{eqn:2.5}) is

\[
\begin{split}
&\Delta_{\sigma}(m)^{-1}\sum_{t \pmod{A\cap K}}
\delta (gt^{-1}\sigma (t))\Phi (gt^{-1}\sigma (t))\\
&=|\gamma_3/\gamma_1|^{-1}\Delta_{\sigma}(m)^{-1}
\sum_{\substack{\Lambda ' =\Lambda '(k,0,k)\in X^*(^{L}A) \\ k\in\zzz}}
\alpha (\Lambda +\Lambda '),\ \ \Lambda =\Lambda (g),
\end{split}
\]
since
$\Delta_G(gt^{-1}\sigma (t))^{-1}\Delta (gt^{-1}\sigma (t))
=|\gamma_3/\gamma_1|^{-1}$. By the previous lemma, the left side of
(\ref{eqn:2.5}) is

\[
\begin{split}
&\Delta_{\sigma}(m)^{-1}|\gamma_3/\gamma_1|^{-1}\eta_1 (\mu )\\
&=\Delta_{\sigma}(m)^{-1}|\gamma_3/\gamma_1|^{-1}
|1-\beta_1/\beta_2||\beta_2/\beta_1|^{1/2}
\int_{A_1(F)\setminus H_1(F)} f_1(x^{-1}hx)\ d\ovx ,
\end{split}
\]
where $h = 
\left( \begin{array}{cc} 
\gamma_1/\gamma_3&\ \\
 \ &\gamma_3/\gamma_1\end{array}
\right) \in N_1(g).$
Since
$\beta_1/\beta_2 = \gamma_1^2/\gamma_3^2$, these equations combine to give, for
$g\in A(F)$ and $h\in N_1(g)\cap H_1(F)$,

\begin{equation}
\label{eqn:2.12}
 \int_{G(g\sigma ,F)\setminus G(F)} \phi (x^{-1}g\sigma (x))\ d\ovx
 =\int_{H_1(h,F)\setminus H_1(F)}
(\lambda_1^*\phi )(x^{-1}hx)\ d\ovx .
%\tag 2.12
\end{equation}
This proves the $SL(2)$--fundamental lemma
in the split case since the stable conjugacy class of h contains only
one element.
 
For the $SO(3)$--fundamental lemma, we have $h' = \left(
\begin{array}{ccc}  \gamma_1/\gamma_3&\ &\ \\
 \ &1&\ \\
 \ &\ &\gamma_3/\gamma_1\end{array}\right) \in N_1(g)$.
(We identify $SO(3)$ with $PGL(2)$ via (\ref{eqn:1.3}).)
Due to the symmetry under the Weyl group, we can replace the sum

\[
\sum_{\substack{M=M(k,0,k)\\k\in \zzz }}
\alpha (\Lambda +M)
=\sum_{\substack{ M=M(0,k,0)\\k\in \zzz }}
\alpha (\Lambda +M),
\]
by a sum over the end coordinate: $(0,0,k)$. By Lemma \ref{lemma:2.11}, 
this sum satisfies

\[
\beta_2 (\nu )=
\sum_{\substack{M=M(0,0,k)\\k\in \zzz }}
\alpha (\Lambda +M).
\]
A simple computation verifies that $\Delta_{\sigma }(m)^{-1}
|\gamma_3/\gamma_1|^{-1}|1-\beta_1||\beta_1|^{-1/2}
=|(1+\beta_1)(1+\beta_2)|^{-1/2}.$
From this, (\ref{eqn:2.9}), and the above expression for $h'$, we have

\begin{equation}
\label{eqn:2.13}
\Phi (g,\phi )=|(1+\beta_1)(1+\beta_2)|^{-1/2}
\Phi (h',\lambda_2^*\phi ),
%\tag 2.13
\end{equation}
for $h'\in N_2(g)\cap H_2(F)$. This proves the $SO(3)$--fundamental
lemma in the split case.
% === END ORIGINAL adjoint_lift2.tex ===
% === BEGIN ORIGINAL adjoint_lift3.tex ===


\section{The reduction in the non--split case to buildings}
\label{sec:3}

\subsection{Summary}
\label{sec:3.1}

R. Langlands \cite{La1}, \cite{La3} has shown how to exploit buildings
to prove fundamental lemmas (see also R. Kottwitz \cite{Ko1}, \cite{Ko2}). 
Following \cite{La1}, this section is concerned 
with  reducing the individual terms of the identities

\begin{equation}
\label{eqn:3.1}
\Phi^{T_1,1}(h,\lambda_1^*(\phi ),\omega_h) = 
\Phi^{T,1}(g,\phi ,\omega_g),
%\tag 3.1
\end{equation}
and

\begin{equation}
\label{eqn:3.2}
\Phi(h',\lambda_2^*(\phi ),\omega_{h'}) = 
\pm |(1+\beta_1)(1+\beta_2)|^{1/2}\Phi^{T,\kappa}(g,\phi ,\omega_g),
%\tag 3.2
\end{equation}
down to finite sums, in case $\phi$ is the characteristic function of a
double coset $KtK$, where $t\in A(F)$ and $K:=G(\off )$. 
We will only be interested in the case $t=1$ but the simple reduction
below is just as easy to carry out for general $t$.
For background, some references are \cite{La3} and \cite{Ko1}.

{\bf Notation}

Recall $H_1:=SL(2)$ and $H_2:=PGL(2)$. Let

\[
Q(t) = Q(\Lambda )=Q(k_1,k_2,k_3) := KtK,
\]
where $T\in A(F)$ is represented by a diagonal matrix with entries
$t_i \in F^{\times}\ |t_i|=q^{k_i}$, and $\Lambda\in X^*({^{L}}A)$
is the character associated to the triple $(k_1,k_2,k_3)$ by duality. 
We assume that the entries have been
ordered in such a way that $k_1\geq k_2\geq k_3$. Of course, the triple
$(k_1,k_2,k_3)\in \zzz^3$ is only well--defined up to translation by
$(n,n,n), \ n\in \zzz$. Let

\[
\begin{split}
\phi_{\Lambda} 
&:= (meas\ Q(\Lambda ))^{-1}char\ Q(\Lambda ),\\
f_{1,\Lambda} 
&:= \lambda_1^*(\phi_{\Lambda}),\\
f_{2,\Lambda} 
&:= \lambda_2^*(\phi_{\Lambda}),\\
a(\Lambda ) 
&:=\Phi^{T_1,1}(h,f_{1,\Lambda}),\\
b(\Lambda ) 
&:=\Phi(h',f_{2,\Lambda}),\\
A(\Lambda ) 
&:=\Phi^{T,1}(g,\phi_{\Lambda}),\\
B(\Lambda ) 
&:=\pm
|(1+\beta_1)(1+\beta_2)|^{1/2}\Phi^{T,\kappa}(g,\phi_{\Lambda}),\\
r(\Lambda )
&:= meas(H_{1}(h,F)\cap H_1(\off ))Qq^{\lambda_1-\lambda_3}a(\Lambda ),\\
s(\Lambda )
&:= meas(H_{2}(h',F)^o\cap H_2(\off ))Qq^{\lambda_1-\lambda_3}b(\Lambda ),\\
R(\Lambda )
&:= meas(G(g\sigma,F)\cap G(\off ))Qq^{\lambda_1-\lambda_3}A(\Lambda ),\\
S(\Lambda )
&:= meas(G(g\sigma,F)\cap G(\off ))Qq^{\lambda_1-\lambda_3}B(\Lambda ),
\end{split}
\]
where
$Q:=(1+q^{-1})(1+q^{-1}+q^{-2})$ and where the $\pm$ sign will be
explained in (\ref{eqn:3.18}) below.
In terms of the above notation, the object of section \ref{sec:5} 
of this paper is to prove that $S(0) = s(0)$.

\subsection{The reduction}
\label{sec:3.2}

To reduce $a(\Lambda )$ and $b(\Lambda )$ down to finite sums, one may
use Macdonald's formula \cite{M} and Plancherel's formula for the Satake
transform, following \cite{La1}. The only special case of this we shall
need is the following

\begin{lemma}
\label{lemma:3.3} 
Suppose that both $h\in H_1(F)$ and $h'\in H_2(F)$
are elliptic regular. In the notation above, we have

\[
r(0)=
\begin{cases} 
Q|{\rm Fix}(h)|,&\ \ {\rm if}\ E=F(\beta_1) \ {\rm is\ unramified},\\
2^{-1}Q|{\rm Fix} (h)|,&\ \ {\rm if}\ E\ {\rm is\ ramified},
\end{cases} 
\]
and

\[
s(0)=
\begin{cases} 
r(0),&\ \ { if }\ \Delta (h)=\Delta (h'),\\
0,&\ \ {\rm if}\  h'\ {\rm has\ no\ fixed\ points},\\
Q,& \ \ {\rm otherwise}.
\end{cases} 
\]
\end{lemma}
 
I shall indicate two different proofs of the first part of 
Lemma \ref{lemma:3.3}
(the first is Lemma \ref{lemma:3.4} (b) and the second is 
(\ref{eqn:3.12})-(\ref{eqn:3.14}) below).
 
The proof
of the following lemma, due to R. Langlands, provides an interesting
application of the Macdonald and Weyl character formulas to the computation of
the number of fixed points in the unramified case.

\begin{lemma}
\label{lemma:3.4} 
First, assume that $E/F$ is unramified, 
where $E:=F[h]=F[h']$ denotes the splitting
field of the torus $T_1:=H_{1,h}$ determined by h. 

\begin{itemize}

\item[(a)] 
If $f_1\in \hhh (SL(2)),\ h\in GL(2,\off )$ then

\[
\Phi^{T_1,1}(h,f_1) = {1\over{meas\ T(\off )}}\int_{^{L}S(\ccc )}
\hat{f_1}(s){q+1\over{q-1}}\big[ {\Delta (h)^{-1}\over{q}}
-{1\over{q}}|c(s)|^{-2}\big]\ ds,
\]
where the notation ($T_1$, $^LS$, $c(s)$) is explained below.

\item[(b)] 
If $f_1 =f_{1,(0,0,0)}$ is the characteristic function of $SL(2,\off )$
divided by the measure of $SL(2,\off )$ then

\[
\Phi^{T_1,1}(h,f_1)=Q|{\rm Fix} (h)|.
\]
\end{itemize}
\end{lemma}

{\bf Notation}
Here $T_1$ denotes the centralizer of $h$ in $H_1$ and
$^{L}S(\ccc )$ is the maximal compact subgroup of $^{L}H_1(\ccc )$
consisting of all $s\in {^{L}H_1}(\ccc )$ with eigenvalues a,b having the
same absolute value. The Haar measure of $^{L}S(\ccc )$ is normalized so 
that its total volume is $1$. Also, here

\[
c(s) := {1-q^{-1}ab^{-1}\over{1-ab^{-1}}}
\]
is the $p$-adic analog of the Harish--Chandra $c$--function for $SL(2)$
(\cite{M}, p. 51), so that the Plancherel measure is given by

\[
{1+q^{-1}\over{2}}|c(s)|^{-2}ds
\]
(see \cite{M}, p. 65, \cite{La3}, p. 46).

\pf (of Lemma \ref{lemma:3.4})

(a) Substitute Lemmas 5.3 and 5.5 into Lemma 5.6 of \cite{La3},
observing that $\Phi^{T_1,1}(h,f_1)=\Delta (h)^{-1}F(h,f_1)$
(in the notation of \cite{La3}, p. 48).

(b) We want to calculate the Satake transform of ${f_{1,\Lambda}},
\ \Lambda =(k_1,k_2,k_3)\in\zzz^3, \ k_1\geq k_2\geq k_3$, in the
case where $\Lambda =(0,0,0)$. By Macdonald's formula 
\cite{M}, p. 52,

\begin{equation}
\label{eqn:3.5}
\hat{f_{1,\Lambda}}(s)
=\sss \big(s_1^{k_1}s_2^{k_2}s_3^{k_3}\prod_{\alpha >0}
{1-q^{-1}\alpha (s)^{-1}\over{1-\alpha (s)^{-1}}}\big)
\cdot Q^{-1}q^{k_3-k_1},
%\tag 3.5
\end{equation}
where the product runs over all positive roots of $sl(3)$ and
$\sss$ is the {\bf symmetry operator} on polynomials in the 
$s_i$ defined by

\begin{equation}
\label{eqn:3.6}
\sss (P(s_1,s_2,s_3)) := \sum_{\sigma\in S_3}
P(s_{\sigma (1)},s_{\sigma (2)},s_{\sigma (3)}),
%\tag 3.6
\end{equation}
where $S_3$ denotes the symmetric group on three letters. Let
${\cal A}$ denote the {\bf anti--symmetry operator} defined by

\begin{equation}
\label{eqn:3.7}
{\cal A} (P(s_1,s_2,s_3)) := \sum_{\sigma\in S_3}
(sgn\  \sigma )P(s_{\sigma (1)},s_{\sigma (2)},s_{\sigma (3)}),
%\tag 3.7
\end{equation}
so that

\begin{equation}
\label{eqn:3.8}
{\cal A} (\sigma P) = (sgn\ \sigma ){\cal A} (P),\ \ \ \sigma\in S_3.
%\tag 3.8
\end{equation}
Let $\rho$ denote half the sum of the positive roots.
Observe that the {\bf Weyl function},

\begin{equation}
\label{eqn:3.9}
\begin{split}
q(s)
&:=\rho (s)\prod_{\alpha >0}(1-\alpha (s)^{-1})\\
&=s_1s_3^{-1}(s_1^{-1}s_2)(1-s_2^{-1}s_3)(1-s_1^{-1}s_3),
\end{split}
%\tag 3.9
\end{equation}
is anti--symmetric in the sense of (\ref{eqn:3.8}) and, by the Weyl character
formula,

\begin{equation}
\label{eqn:3.10}
\begin{split}
\sss \big(\prod_{\alpha >0}
{1-q^{-1}\alpha (s)^{-1}\over{1-\alpha (s)^{-1}}}\big)
&={{{\cal A} \big(\rho (s)\prod_{\alpha >0}(1-q^{-1}\alpha (s)^{-1})\big)}
\over{\rho (s)\prod_{\alpha >0}(1-\alpha (s)^{-1})}}\\
&=Q\cdot {\rm tr}\, (r_{trivial}) = Q,
\end{split}
%\tag 3.10
\end{equation}
where $r_{trivial}$ is the trivial representation on $^{L}H_1$.
(In fact, for arbitrary $\Lambda$, $\hat{f_{\Lambda}}$ can be expressed
as a linear combination of characters of finite dimensional representations 
of $^{L}H$.) The result now
follows from (\ref{eqn:3.12}) (or Lemma \ref{lemma:4.5}) below and
Plancherel's formula \cite{La3}, Lemma 5.4.
\qed

 
For the moment, $E/F$ need not be unramified. 
Let $h\in SL(2,F)$ be non--split,
so $E/F$ is quadratic and $T:=H_1(h)=H_1(h)^o$ is a non--split Cartan. 
We have

\begin{equation}
\label{eqn:3.11}
\begin{split}
\int_{T(F)\setminus H_1(F)} f_{1,(0,0,0)}(x^{-1}hx)\ d\ovx
&={1\over{meas\ T(F)}}
\int_{H_1(F)} f_{1,(0,0,0)}(x^{-1}hx)\ d\ovx\\
&={1\over{meas\ T(F)}}|{\rm Fix} _0(h)|,
\end{split}
%\tag 3.11
\end{equation}
where, in the notation of \cite{Ko2},

\[
{\rm Fix}_0(h):=\{x\in X_{H_1}(0)\ |\ hx=x\}.
\]
If $h^\delta\in \ddd (h),\ T^\delta$ the associated Cartan, for some
$\delta\in GL(2,F)$, then

\[
\int_{T^\delta (F)\setminus H_1(F)} f_{1,(0,0,0)}(x^{-1}h^\delta x)\ d\ovx
= {1\over{meas\ T(F)}}|{\rm Fix} _0(h^\delta )|,
\]
since the Haar measures on $T$ and $T^\delta $ have been chosen so that
$meas\ T(F)=meas\ T^\delta (F)$ (we have $T^\delta \cong T$ over $\ovF$ and
the Haar measure on $T^\delta $ is defined by the pull--back of that 
on $T$).
 
The following facts (which can be found in \cite{Bl}, ch. 1), 
play a role here: if $h$ splits over a
ramified extension, then the fixed points of $h$ (or of $h^\delta $) in
the tree of $H_1(F)$ are those less than a certain distance
from a certain point in the first barycentric
subdivision of $X_{H_1(F)}$; if $h$ splits over an unramified extension
then the fixed points are those less than a certain distance from a 
certain vertex in $X_{H_1(F)}$. In the ramified case the fixed 
points occur in pairs, one in $X(0)$ and
the other in $X(1)$. In particular, if $T$ is ramified then 
$|{\rm Fix} _0(h)| = |{\rm Fix} _0(h^\delta )|$ and
$|{\rm Fix} _1(h)| = |{\rm Fix} _1(h^\delta )|$. Writing

\[
{\rm Fix}(h) = {\rm Fix} _0(h)\cup {\rm Fix}_1(h),
\ \ {\rm (disjoint\ union)},
\]
where ${\rm Fix}_i(h):= X_{H_1(F)}(i)\cap {\rm Fix} (h)$, we have
in the unramified case

\[
|{\rm Fix} _0(h)|=|{\rm Fix} _1(h^\delta )|,
\ \ \ |{\rm Fix} _0(h^\delta )|=|{\rm Fix} _1(h)|,
\]
and therefore, by (\ref{eqn:1.23}),

\begin{equation}
\label{eqn:3.12}
\Phi^{T,1}(h,f_{1,(0,0,0)}) = {1\over{meas\ T(F)}}|{\rm Fix} (h)|.
%\tag 3.12
\end{equation}
A similar result holds for $\Phi_{H_2}(h',f_{2,(0,0,0)})$, where
$h'\in N_2(g)\cap H_2(F)$.
 
Suppose $|{\rm det }\, h| =|\pi |^{2n}$, for some $n\in\zzz$. This is 
necessary for Fix$(h)\not=\emptyset$. If $E/F$ is ramified then the
orbits on $X_{H_1}$ under $T(F)$ are twice as large as the orbits under
$T(F)\cap H_1(\off )$, so that

\begin{equation}
\label{eqn:3.13}
meas\ T(F) = 2\cdot meas(T(F)\cap H_1(\off ))
%\tag 3.13
\end{equation}
(see \cite{La3}, p. 52); if $E/F$ is unramified then

\begin{equation}
\label{eqn:3.14}
meas\ T(F)= meas(T(F)\cap H_1(\off )).
%\tag 3.14
\end{equation}
These follow from the classification of the Cartan 
subgroups of $H_1(F)$ in terms
of the extensions of $F$. Recall that Haar measures have already been choosen 
so that 

\[
meas(G(g\sigma ,F)\cap G(\off ))=meas(H_1(h,F)\cap H_1(\off )).
\]
This proves the first part of (\ref{lemma:3.3}). 
The second part of (\ref{lemma:3.3}) follows along 
similar lines.
 
Langlands' proof that $R(\Lambda )=r(\Lambda )$ uses the identity (actually
a facsimile thereof, given below)

\begin{equation}
\label{eqn:3.15}
R(\Lambda )=Q(\Lambda )q^{k_3-k_1}\sum_{g'\in \ddd (g)}
\sum_{\substack{P\\Inv(g'\sigma (P),P)=\Lambda}} 1,
%\tag 3.15
\end{equation}
where

\[
Q(\Lambda )=
\begin{cases} 1,&\ \ \ \ \ \ k_1>k_2>k_3,\\
1+q^{-1},&\ \ \ k_1=k_2, {\rm  or  }\, k_2=k_3,\ {\rm  but }\ k_1\not= k_3,\\
Q,&\ \ \ \ \ \ \ \Lambda = (0,0,0),
\end{cases} ,
\]
and the inner sum runs over the vertices in the building for
$G(F)$ (the inner sum is independent of the representative $g'$ choosen 
for the $\sigma$--conjugacy class). In case $N_1(g)$ has distinct
eigenvalues not in $F$, by explicitly analyzing the proof of Lemma 
\ref{lemma:1.21} one
can show that, as sets,

\begin{equation}
\label{eqn:3.16}
\ddd_{\sigma}(g) \cong H^1(F,G(g\sigma )) 
\cong F^{\times}/N_{_{E/F}}(E^{\times}),
%\tag 3.16
\end{equation}
where $E$ denotes the quadratic extension determined 
by the eigenvalues of $N_1(g)$. 
Very briefly, one uses (\ref{eqn:3.16}) 
to parameterize $\ddd_{\sigma}(g)$
in terms of certain representatives $g_a,\ a\in F^{\times}$. 
One can
show that this parameterization satisfies the following property:
two such $g_a,\ g_b$ belong to the same class in $\ddd_\sigma (g)$ 
if and only if
$a\cdot N_{_{E/F}}(E^{\times})=b\cdot N_{_{E/F}}(E^{\times})$. 
Rather than (\ref{eqn:3.15}), it is actually 
necessary to use the more general identity

\begin{equation}
\label{eqn:3.17}
R(\Lambda )={2\over{[F^{\times}:U]}}
Q(\Lambda )q^{k_3-k_1}\sum_{a\in F^{\times}/U}
\sum_{\substack{ P\\Inv(g_a\sigma (P),P)=\Lambda}} 1,
%\tag 3.17
\end{equation}
where $U\subseteq N_{_{E/F}}(E^{\times})\subseteq F^{\times}$
is a compact open subgroup of finite index. Similarly, we have

\begin{equation}
\label{eqn:3.18}
S(\Lambda )={\pm 2\over{[F^{\times}:U]}}|(1+\beta_1)(1+\beta_2)|^{1/2}
Q(\Lambda )q^{k_3-k_1}\sum_{a\in F^{\times}/U}
\sum_{\substack{P\\Inv(g_a\sigma (P),P)=\Lambda}} 
\kappa (g_a),
%\tag 3.18
\end{equation}
where $\kappa :\ddd_{\sigma }(g) \rightarrow \{\pm 1\}$
is defined to be $+1$ on the element of $H^1(F,G(g\sigma ))$
corresponding to the split inner form of $G(g\sigma ))$ and to be
-1 on the other element. The $\pm$ sign is choosen so that 
$\pm\kappa(g_a)=(a,\epsilon)_2$, where $E=F(\sqrt{\epsilon})$
and $(\ ,\ )_2$ denotes the quadratic Hilbert symbol
attached to $F$. The sign $\kappa$ satisfies the following property:

\begin{lemma}
\label{lemma:3.19}
%{Fact 3.19} 
For any $g',\ g'' \in G(F)$, we have
$\kappa (g')=\kappa (g'')$ if and only if $g'$ and $g''$ belong to the
same class in $\ddd_\sigma (g)$.
\end{lemma}
 
The expressions (\ref{eqn:3.17}), (\ref{eqn:3.18}), 
and those in Lemma \ref{lemma:3.3} constitute 
the desired reductions.

% === END ORIGINAL adjoint_lift3.tex ===
% === BEGIN ORIGINAL adjoint_lift4.tex ===

\section{The buildings for $PGL(3)$ and $SL(2)$}
\label{sec:4}

For background, see for example \cite{Ko1} and \cite{T}.
The results (and pictures) in this section are, as 
mentioned earlier, due to R. Langlands
\cite{La1}. We assume $p\not= 2$ throughout this section.
 
Fix a vector space $V$ over $F$ of dimension either 2 or 3.
We may consider the buildings $\bbb (G),\ \bbb (H_1)$ for $G,\ H_1$
as graphs of lattice classes; if L is a rank two or three $\off$--lattice
in $V$
then its class is denoted by [L]. Recall that two vertices $[L_1],\ [L_2]$
are to be joined by an {\bf edge} if and only if there exist 
$\lambda_1,\ \lambda_2 \in F^{\times}$ such that

\begin{equation}
\label{eqn:4.1}
\pi\lambda_1 L_1 \sne \lambda_2 L_2 \sne \lambda_1 L_1.
%\tag 4.1
\end{equation}
The edge from $[L_2]$ to $[L_1]$ is {\bf positively directed} if (4.1)
holds and if $\lambda_1 L_1/\lambda_2 L_2$ is a module of rank one
over the residue field.
Three vertices $L_1,\ L_2,\ L_3$ are the vertices of a 2-simplex
or {\bf chamber} of $\bbb (G)$ if and only if there exists
$\lambda_1,\ \lambda_2,\ \lambda_3 \in F^{\times}$ such that

\begin{equation}
\label{eqn:4.2}
\pi\lambda_1 L_1 \sne \lambda_3L_3 
\sne \lambda_2 L_2 \sne \lambda_1 L_1.
%\tag 4.2
\end{equation}
 
To describe the action of $G$ on
$\bbg$ one must first represent $g \in G(F)$ by a matrix
$A\in GL(3,F)$. The point is that the definition of the lattice class
implies that the center acts trivially, so the action of
$g = A\cdot Z(F)$ is well--defined, where $Z(F)$ denotes the
center of $G(F)$. Explicitly, the action of G(F) on $\bbg$ is given by
$g: [L] \longmapsto [AL].$
The action of $H_1$ on $\bbH1$ is simpler: $h\in H_1(F)$ sends
$[\Lambda ]$ to $[h\Lambda ]$. These actions send vertices to vertices, 
edges to edges, chambers to chambers
(in the case of $\bbg$), preserve orientation and are
automorphisms of the buildings. 
 
The action of $\sigma$ on $\bbg$, however, does not preserve orientation.
This ``anti--automorphism'' sends a lattice $L$ to its dual lattice

\begin{equation}
\label{eqn:4.3}
\check{L} :=\{v \in V\ |\ {^{t}v}\cdot v' \in \off , \ \forall v' \in L\}.
%\tag 4.3
\end{equation}
If $L=g\Lambda_0$, where $\Lambda_0$ denotes the lattice class of
$\off^3$, then this means

\[
\sigma : g\Lambda_0 \longmapsto {^{t}g^{-1}}\Lambda_0.
\]
The composition $g\sigma$ acts by sending $[L]$ to $[A\check{L}]$. It is the
fixed point set of $g\sigma$ in $\bbg$ that we want to describe.

\begin{lemma}
\label{lemma:4.4} 
\begin{itemize}

\item[(a)] 
Suppose that the number of fixed points of 
$h\in H_1(F)$ in $\bbH1$ is finite: $|{\rm Fix}(h)| < \infty$. 
(This is the case if
$h$ is elliptic regular, i.e., the eigenvalues of $h$ generate a quadratic
extension of $F$.) Then, for any $x\in H_1(F)$, we have 
$|{\rm Fix}(h)| =|{\rm Fix}(x^{-1}hx)|$,
i.e., the number of fixed points depends only on the conjugacy class of $h$.

\item[(b)] 
Suppose that the number of fixed points of $g\sigma$ in $\bbg$ is
finite: $|{\rm Fix}(g\sigma )|<\infty$. (This is the case if $N_1(g)$ is 
elliptic regular.) Then, for any $x\in PGL(3,F)$, we have 
$|{\rm Fix}(g\sigma )| = |{\rm Fix}((x^{-1}g\sigma (x))\sigma )|$,
i.e., the number of fixed points of $g\sigma$ depends only on the
$\sigma$--conjugacy class of $g$.
\end{itemize}
\end{lemma}

\begin{remark}
The proof of this lemma is easy, hence is omitted. 
If $h,\ h'$ are stably conjugate but not conjugate over $F$ one
may have $|{\rm Fix}(h)|\not= |{\rm Fix}(h')|$, 
and a similar statement holds for
$\sigma$--conjugacy.
\end{remark}
 
The following result is well-known. 

\begin{lemma}
\label{lemma:4.5} 
Let $h\in GL(2,F)$ be an elliptic regular element
with $\Delta_{H_2} (h)=q^{-k}$, for some $k\geq 0$. Let $E/F$ denote the
quadratic extension generated by $h$. 

\begin{itemize}
\item 
If $E/F$ is unramified then

\[
|{\rm Fix}(h)|={{q^{k+1}+q^k-2}\over{q-1}}.
\]        
\item 
If $E/F$ is ramified then

\[
|{\rm Fix}(h)|={2(q^k-1)\over{q-1}}.
\]        
(We have used the fact that the conductor of $E/F$, as defined in \cite{Bl},
is $1$ in the ramified case.)
\end{itemize}
\end{lemma}

We now describe an embedding

\[
\bbb (H_1) \hookrightarrow \bbb (G).
\]
First, replace the tree $\bbb (H_1)$ by the product of itself with 
the affine line:
replace all vertices by lines (copies of $V_A$ running perpendicular to the 
edges, as visualized in three dimensions) and all edges by strips. 
This figure may be regarded as $\bbb (GL(2))$. 
In order to embed this into
$\bbb (G)$, one must impose a simplicial structure on it compatible with that
of $\bbb (G)$. 
Each line $\overline{\Lambda}$ of $\bbb (GL(2))$ associated to a
vertex $\Lambda$ of $\bbb (H_1)$ must first be provided with its 
own vertices. 
We define the vertices on $\overline{\Lambda}$ to be 
$\Lambda_{(n)} :=\pi^{-n}\off + \Lambda$,
for fixed embeddings
$\pi^{-n}\off \hookrightarrow V_A,\ \Lambda\hookrightarrow W_A$. 
Here $V_A, \ W_A$ are as in subsection 1.3. Given two neighboring
vertices $\Lambda , \Lambda '$ of $\bbb (H_1)$ we obtain neighboring lines
$\overline{\Lambda},\ \overline{\Lambda '}$ and the vertices of these lines 
are joined by an edge if they are neighbors. 
This simplicial structure 
on $\bbb (GL(2))$ defines an embedding of
$\bbb (H_1)$ into $\bbb (G)$ whose image we denote by $\zz = \zz (V_A,W_A)$. 
Let

\[
\pi\lambda \Lambda \sne \lambda ' \Lambda ' 
\sne \lambda \Lambda ,
\]
and let $L'_{(n)} := \pi^{-n}\off + \Lambda ', 
\ L_{(n)}:=\pi^{-n}\off + \Lambda$. The band in $\zz$ associated to the edge 
joining $\Lambda $ and $\Lambda '$ in $\bbH1$ is represented by a 
picture of the following sort:

%%figure (4.6) on page 1119
\begin{figure}[!h!]
\begin{picture}(200.00,100.00)(-120.00,0.00)
\thicklines
\put(-80.00,50.00){\line(1,0){220.00}} 
\put(-10.00,60.00){$L_{(-1)}$} 
\put(-2.00,47.00){$\bullet$} 
\put(0.00,50.00){\line(1,-1){50.00}}
\put(100.00,50.00){\line(1,-1){30.00}}
\put(-60.00,-23.00){$L'_{(-1)}$} 
\put(-52.00,-3.00){$\bullet$} 
\put(0.00,50.00){\line(-1,-1){50.00}}
\put(47.00,-3.00){$\bullet$}  
\put(90.00,60.00){$L_{(0)}$} 
\put(-80.00,0.00){\line(1,0){220.00}}
\put(40.00,-23.00){$L'_{(0)}$} 
\put(50.00,0.00){\line(1,1){50.00}}
\put(-50.00,0.00){\line(-1,1){20.00}}
\put(98.00,47.00){$\bullet$}  
\put(158.00,43.00){$\Lambda$}  
\put(158.00,-7.00){$\Lambda'$}  
\end{picture}
\vskip .3in
\caption{The band in the apartment \label{fig:4.6}}
\end{figure}
\vskip .4in

\noindent
The lines in $\zz$ have been drawn in such a way that the positive direction 
on them is from left to right. As mentioned earlier, the action of
$\text{Aut}_F(W_A) \hookrightarrow \text{Aut}_F(W_A +V_A)$ on $\zz$ preserves
orientation. 
The action of $\sigma$ on $\bbg$ does not even restrict to an action
on $\zz = \zz (V_A,W_A)$. However, the action of $g\sigma$, with 
$g=[A] \in G(F)$ 
represented by $A \in GL(3,F)$, does send $\zz (V_A, W_A)$ to
$\zz (AW_A^{\bot}, AV_A^{\bot}) = \zz (V_A, W_A)$, reversing orientation. 
Here we have used the fact the $AW_A^{\bot} = V_A$ and 
$AV_A^{\bot}=({^{t}A^{-1}}V_A)^{\bot}$, by definition of $W_A$.
As we shall see, thanks to Corollary \ref{cor:1.14} it turns out 
to be sufficient
for our purpose of relating fixed points of $g\sigma$ 
in $\bbg$ with fixed points
of $h\in N_1(g),\ h\in H_1(F),$ in $\bbH1$ to assume that the matrix A in
$GL(3,F) \cong \text{Aut}_F(V_A +W_A)$ is in $(1,2)$-block form. 
For such
matrices, it is possible to explicitly realize the action of $g\sigma$
on $\zz$ in terms of the geometry described above. 
In fact, this is what we do next.
 
In \cite{La1}, R. Langlands introduced the notion of a ``characteristic
leaf'' to understand explicitly vertices in $\bbg$ not in $\zz$ in terms
of vertices in $\zz$. 
This notion is very useful for relating fixed
points of $g\sigma$ in $\bbg$ to fixed points of $h\in N_1(g)$ in $\bbH1$. 
Given
$L:=\pi^n\off +\pi^a\off +\pi^b\off$, let $M_1:=\pi^n\off,\ M_2:=\pi^b\off,
\ M_3:=\pi^{-b}\off$. 
Consider the vertices of the {\bf characteristic
leaf} for $L$:

\begin{equation}
\label{eqn:4.7}
\begin{split}
&L':= M_1 +M_2 +M_3 \in \zz,\\ 
&L:= \pi^b M_1 + M_2 +M_3 \notin \zz,\\ 
&L'':= \pi^b M_1 +\pi^b M_2 + M_3 \in \zz,
\end{split}
%\tag 4.7
\end{equation}
depicted in Figure \ref{fig:4.8}.

\begin{figure}[!h!]
\begin{picture}(200.00,150.00)(-120.00,0.00)
\thicklines
\put(-50.00,100.00){\line(1,0){220.00}}
\put(-50.00,0.00){\line(1,0){220.00}}
\put(71.00,50.00){$\bullet$} 
\put(80.00,-20.00){$L''$} 
\put(86.00,-2.00){$\bullet$} 
\put(135.00,97.00){$\bullet$} 
\put(60.00,50.00){$L$} 
\put(140.00,110.00){$L'$} 
\put(90.00,00.00){\line(3,6){50.00}}

\thinlines
\put(90.00,00.00){\line(-1,3){17.00}}
\put(72.00,52.00){\line(4,3){65.00}}

\end{picture}
\vskip .3in
\caption{A characteristic leaf \label{fig:4.8}}
\end{figure}

\vskip .3in
\noindent
The vertices of the simplices forming the equilateral triangle
are given by
$\pi^{n_1}M_1 +\pi^{n_2}M_2 + M_3,\ 0 \leq n_2 \leq n_1 \leq b$.
The segment of length $|b|$ joining $L'$ to $L''$ is called the
{\bf characteristic base} or {\bf segment}. This base degenerates to a
vertex if $L \in \zz$ and otherwise it may be regarded as 
an equilateral triangle, under a suitable metric structure on $\bbg$.

\begin{lemma}
\label{lemma:4.9} 
(\cite{La1}) 
Suppose $g\in G(F)$ is such that $N_1(g)$
is elliptic regular. Then the following hold.

\begin{itemize}
\item[(a)] 
$g\sigma$ sends the characteristic leaf of $L$ to the characteristic leaf
of $g\check{L}$.

\item[(b)] 
For each vertex $\Lambda$ in $\bbH1$, $g\sigma$ sends the line 
$\{\pi^{-n}\off +\Lambda\ |\ n\in \zzz\}$
of $\zz$ with positive orientation to another line 
$\{\pi^m\off +\Lambda ' \ |\ m \in \zzz\}$ in
$\zz$ with reversed orientation.

\item[(c)] 
$g\sigma$ sends the characteristic base of the leaf for L to the
characteristic base (with orientation reversed) for $g\check{L}$;

\item[(d)] 
For any $h\in SL(2,F)$ representing $N_1(g)$, the action of
$h$ on $\zz$ is the same as the action of $(g\sigma)^2$ restricted
to $\zz$.
\end{itemize}
\end{lemma}

\begin{remark}
The $h \in H_1(F)$ in (d) is fixed for the rest of this section.
\end{remark}

\pf
We may assume by Corollary \ref{cor:1.14} that $g$ is represented by a 
$(2,1)$-block matrix $A\in GL(3,F)$ in the basis giving the embedding
$\zz$. From this the four statements (a)-(d) follow without difficulty
from simple facts about buildings. 
\qed

In particular, (d) implies that if $\overline{\Lambda '},
\ \overline{\Lambda ''}$ are lines associated to two fixed points
$\Lambda ', \Lambda ''$ of $h \in N_1(g)$ and if the endpoints $L'$, $L''$
of the vertices of the characteristic base lie on $\overline{\Lambda '},
\ \overline{\Lambda ''}$, respectively, then $g\sigma$ ``flips'' $L$
and its leaf about a certain ``dotted'' line $D_A$ (which depends only
on $A$, not on $L$, $L'$, $L''$, and lies in the first barycentric
subdivision of $\bbg$ but does not necessarily lie in $\bbg$):


%%figure (4.10) on page 1121
\begin{figure}[!h!]
\begin{picture}(200.00,150.00)(-120.00,0.00)
\thicklines
\put(-50.00,100.00){\line(1,0){250.00}}
\put(-50.00,0.00){\line(1,0){250.00}}
%left leaf
\put(-29.00,50.00){$\bullet$} 
\put(-20.00,-20.00){$L''$} 
\put(-14.00,-2.00){$\bullet$} 
\put(35.00,97.00){$\bullet$} 
\put(-35.00,33.00){$L$} 
\put(40.00,110.00){$L'$} 
\put(-10.00,00.00){\line(3,6){50.00}}

%right leaf
\put(91.00,50.00){$\bullet$} 
\put(100.00,-20.00){$M''$} 
\put(106.00,-2.00){$\bullet$} 
\put(155.00,97.00){$\bullet$} 
\put(86.00,33.00){$M$} 
\put(160.00,110.00){$M'$} 
\put(110.00,00.00){\line(3,6){50.00}}

\thinlines
%right leaf
\put(110.00,00.00){\line(-1,3){17.00}}
\put(92.00,52.00){\line(4,3){65.00}}

%left leaf
\put(-10.00,00.00){\line(-1,3){17.00}}
\put(-28.00,52.00){\line(4,3){65.00}}

%middle line
\put(50.00,00.00){\line(3,6){50.00}}
\put(43.00,33.00){$D_A$} 

\put(-110.00,0.00){$\Lambda'' =\gamma \Lambda''$} 
\put(-110.00,100.00){$\Lambda' =\gamma \Lambda'$} 
\end{picture}
\vskip .3in
\caption{Flipping characteristic leafs \label{fig:4.10}}
\end{figure}

\vskip .4in
\noindent
If $N_1(g)$ does not contain a unipotent then by Corollary 
\ref{cor:1.14} we may assume that the
matrix $A\in GL(3,F)\cong \text{Aut}_F(V_{A'}+W_{A'})$ 
representing $g=[A]$ is in block form:

\begin{equation}
\label{eqn:4.11}
A = 
\left(
\begin{array}{cc}
u\pi^c&\ \\\ &C
\end{array}
\right) ,
\ \ \ C\in GL(2,F),\ u\in \off^{\times},\ c\in\zzz . 
%\tag 4.11
\end{equation}
(Here $[A']$ is the matrix in Corollary \ref{cor:1.14} 
which belongs to the same stable
$\sigma$--conjugacy class as $[A]$.)
In this case, $A\sigma$ sends the lattice $\pi^{-n}\off +\Lambda$ in $\zz$ to
$\pi^{c+n}\off +\Lambda '$, for some $\Lambda '$. If we parameterize
the vertices in the lines in Figure \ref{fig:4.8} 
by the integers then, in a sense, 
$D_A$ may be thought of 
as the ``$c/2$--line''. It belongs to $\bbg$ if and only if $c$ is even. 
Suppose that $c$ is even and that 

\begin{equation}
\label{eqn:4.12}
\begin{split}
&A\sigma : [\pi^{-n}\off +\Lambda] \longmapsto [\pi^{c+n}\off +\Lambda '],\\
&A\sigma : [\pi^{-n}\off +\Lambda '] \longmapsto [\pi^{c+n}\off +\Lambda].
\end{split}
%\tag 4.12
\end{equation}
It can be shown without much effort that, in this case,  
$[\Lambda ],\ [\Lambda ']$ in $\zz$ correspond to two fixed points of
$h \in N_1(g)$ in $\bbH1$. Also, it is not hard to see that,

\begin{lemma}
\label{lemma:4.13}
%{Fact 4.13} 
Suppose c is even and L is a fixed point of
$[A]\sigma$ whose characteristic base has endpoints lying on the
lines $\overline{\Lambda '}\not= \overline{\Lambda ''}$. Then 
the following hold.

\begin{itemize}
\item[(a)] 
$D_A$ intersects the lines 
$\overline{\Lambda '},\ \overline{\Lambda ''}$
at the two vertices which determine the endpoints of the characteristic base 
of L and

\item[(b)] 
The lines $\overline{\Lambda '},\ \overline{\Lambda ''}$ in $\zz$
correspond to two fixed points of $h$ in $\bbH1$.
\end{itemize}
\end{lemma}

\begin{remark} 
There is a simpler version of this dealing with the case 
$\overline{\Lambda '}= \overline{\Lambda ''}$.
In this way, we will obtain a ``2-1'' correspondence, when
$c$ is even, between \underbar{pairs} of fixed points of $h$ in $\bbH1$ and
fixed points of $[A]\sigma$ in $\bbg$. One important point to bear in mind is 
that when $c$ is \underbar{odd} then the line $D_A$ cannot intersect $\bbg$
and therefore there can be \underbar{no} fixed points of $[A]\sigma$ in this
case. So far we have only shown that (when $c$ is even) given a fixed point 
of $[A]$ in $\bbg$ there are associated two fixed points
of $h$ in $\bbH1$.
\end{remark}
 
From the reduction of \S \ref{sec:3} 
one sees that it is also necessary to know what
happens when $A$ in (\ref{eqn:4.11}) is replaced by a matrix $A'$ 
stably $\sigma$--conjugate
to $A$ but not $\sigma$--conjugate to $A$ over $F$. Of course, 
$N_1([A])=N_1([A'])$ so $A'$ may be put in $(1,2)$-block 
form as in (\ref{eqn:4.11}),
for some $c'\in\zzz ,\ C'\in GL(2,F),\ u'\in\off^{\times}$. 
If moreover the eigenvalues of $h$ determine an unramified 
extension of $F$ then
$c'$ and $c$ have different parity. On the other hand, if $A'$ and $A$
are $\sigma$--conjugate over $F$ then $c$ and $c'$ do have the same parity, so
we have the following result.

\begin{lemma}
\label{lemma:4.14}
%{Fact 4.14} 
The parity of $c$ depends only on the stable 
$\sigma$--conjugacy class of $[A]$. If the eigenvalues of $h$ determine an
unramified extension of $F$ (the ``unramified case'') and $A$, $A'$ are stably 
$\sigma$--conjugate but not
$\sigma$--conjugate over $F$ then either $c$ or $c'$ is even (but not both). 
\end{lemma}

By the previous paragraph, in this ``unramified case'', there is either 
a fixed point of $A\sigma$ (and $A'\sigma$ has no fixed point) or a
fixed point of $A'\sigma$ (and $A\sigma$ has no fixed point). 
 
The problem now is to deal with the converse of Lemma \ref{lemma:4.13}: 
given a pair of fixed points
of $h \in N_1(g)$, associate to them fixed point of $A\sigma$. Instead
of proving this, we shall prove that this holds ``on average'', which is
sufficient for our purposes.
For this the following lemma is crutial, but first we need some notation:
Let $d$ be the maximum distance
between two fixed points of $h$, let the eigenvalues of $h$
determine a quadratic extension $E$ of $F$, and suppose that

\begin{equation}
\label{eqn:4.15}
U(d):=U_d\pi^{2\zzz}:=
\{x\in\off^{\times}\ |\ x\equiv 1 \, {\rm mod}\, \pi^d\}\pi^{2\zzz}
\subseteq N_{_{E/F}}(E^{\times})
%\tag 4.15
\end{equation}
is finite index in $F^{\times}$. Let ${\cal S}$ denote the set of 
$a\in F^{\times}$ for which
$P'\not= P''$ exist (for a fixed $p'$, $p''$) and satisfying (ai)-(aiv) in the 
lemma below. Here $A_a$ is defined to be a representative of $g_a$, where 
$g_a$ is defined in the paragraph following (\ref{eqn:3.16}) and $A$ is our
representative of $g$ choosen at the beginning of this section. We may,
by Corollary \ref{cor:1.14}, choose all the $A_a$ to be in block form as in 
(\ref{eqn:4.11}).
$A_a\sigma$ may have no
fixed points; however, if it does, let $P_0$ denote one of them. Let
$U_n$ denote the stabilizer of the action of 
$T_a :=(A_a\sigma )(A\sigma )^{-1}$ on $P_0$, for some $n\geq 0$, with
$U_n$ as in (\ref{eqn:4.15}).

\begin{lemma}
\label{lemma:4.16} 
(\cite{La1}, Lemma 5.1) 
Suppose that $p'$, $p''$ are 
vertices in
$\bbH1$ (possibly $p'$ = $p''$) and that $[A]\sigma$ ``flips'' the lines
in $\bbg$ associated to $p'$, $p''$ onto each other. Then:

\begin{itemize}
\item[(a)] 
For each $a\in F^{\times}$ there exists at most one pair $P'$, $P''$
of vertices in $\bbg$ with

\begin{itemize}
\item[(i)] 
$ P'$ lying on the line associated to $p'$,

\item[(ii)] 
$P''$ lying on the line associated to $p''$,

\item[(iii)]
In the notation of \S \ref{sec:3}, $A_a\sigma (P') = P'',
\ A_a\sigma (P'')=P'$.

\item[(iv)] 
$P'$, $P''$ form the extreme vertices of a characteristic
segment.
\end{itemize}

\item[(b)] 
Suppose $a\in F^{\times}$ is such that a pair $p'$, $p''$ exists for $A_a$.
Then a pair exists for $A_b$ (for the same $p'$, $p''$) if and only
if $aN_{_{E/F}}(E^{\times})=bN_{_{E/F}}(E^{\times})$.

\item[(c)] 
Let ${\cal S},\ P_0,$ and $U_n$ be as above, and assume that $n\geq 1$.
Then $n \leq d$ and there are $[\off^{\times}:U_n]$ distinct
P's with segment $(P',P'')$ such that $(A_a\sigma )P = P$.

\item[(d)] 
If $|{\rm Fix}(A_a\sigma )|\not= 0$ then $|{\rm Fix}(A_{a\pi}\sigma )|=0$,
for $E/F$ unramified.

\item[(e)] 
In the notation of part (c),

\[
{1\over{[\off^{\times}:U_d]}}\sum_{a\in F^{\times}/U(d)}
\sum_{\substack{ P\\A_a\sigma P = P\\P \text{ has segment }(P',P'')}}
1 
={1\over{[\off^{\times}:U_n]}}\sum_{a\in F^{\times}/U(n)}
\sum_{\substack{ P\\A_a\sigma P = P\\P \text{ has segment }(P',P'')}}
1 .
\]
\end{itemize}

There is an obvious ``$\kappa$--analog'' (see (\ref{eqn:3.18})) 
of the above equation 
whose statement is left to the reader.
\end{lemma}

\begin{remark} 
(a) When applying this lemma to (\ref{eqn:3.17}) or 
(\ref{eqn:3.18}), one uses the fact that

\[
{1\over{2}}[F^{\times}:U(n)] = [\off^{\times}:U_n].
\]
(b) For the situation in Lemma \ref{lemma:4.13}(c) when n = 0, 
see \S \ref{sec:5}. The ``matching''
is then $P = P' = P''\longmapsto p'=p''$.
\end{remark}

Let me now sketch the proof of this lemma. Consider the situation
described by the following diagram:

%%figure (4.17) on page 1123
\begin{figure}[!h!]
\begin{picture}(200.00,100.00)(-180.00,0.00)
\thicklines
\put(-100.00,50.00){\line(1,0){270.00}} 
\put(-10.00,60.00){$Q'$} 
\put(-2.00,47.00){$\bullet$} 
\put(0.00,50.00){\line(1,-1){50.00}}
\put(100.00,50.00){\line(1,-1){50.00}}
\put(140.00,-23.00){$Q''$} 
\put(148.00,-3.00){$\bullet$} 
\put(0.00,50.00){\line(-1,-1){30.00}}
\put(47.00,-3.00){$\bullet$}  
\put(90.00,60.00){$P'$} 
\put(-100.00,0.00){\line(1,0){270.00}}
\put(40.00,-23.00){$P''$} 
\put(50.00,0.00){\line(1,1){50.00}}
\put(150.00,0.00){\line(1,1){20.00}}
\put(98.00,47.00){$\bullet$}  
\put(-150.00,20.00){$P',P'',Q',Q''\in {\cal Z}$}  
\end{picture}

\vskip .3in
\caption{Another band in the apartment \label{fig:4.17}}
\end{figure}

\vskip .4in
\noindent
In this case, as we've noted above, there are three rank one $\off$--modules
$M_1,\ M_2,\ M_3$ such that

\[
\begin{split}
&P' = [M_1+M_2+M_3],\ \ P''=[\pi^bM_1+\pi^bM_2 +M_3],\\
&Q' = [\pi^bM_1+M_2+M_3],\ \ Q''=[M_1+\pi^bM_2 +M_3].
\end{split}
\]

\begin{lemma}
\label{lemma:4.18}

\begin{itemize}
\item[(a)] 
The vertices $P$ in $\bbg$ with characteristic segment $(P',P'')$,
$P' \not= P''$, correspond to rank three lattices
$L+M_3$, where $L$ is a lattice satisfying
\newline
$\ \ (i) \pi^bM_1 +\pi^bM_2 \subseteq L \subseteq M_1+M_2,$
\newline
$\ \ (ii) \pi^{b-1}M_1 + \pi^{b-1}M_2 \not\subset L,$
\newline
$\ \ (iii) L \not\subset \pi M_1+\pi M_2,$
\newline
$\ \ (iv)  L \not\subset \pi M_1+ M_2 =:M',$
\newline
$\ \ (v) \  L \not\subset  M_1+\pi M_2=:M''.$
\newline

\item[(b)] 
If $u\in \off^{\times}$ then the action of 
$T_u := (A_u\sigma )(A\sigma )^{-1}$ on $\zz$ induces an action of 
$\off^{\times}$ on the lattices $L$ in (a). This action agrees with
that defined by the matrix

\[
\left(
\begin{array}{cc}
u&0\\0&1 
\end{array}
\right),
\]
and, moreover, it is transitive on the set of $L$'s.
\end{itemize}
\end{lemma}

\pf
(of Lemma \ref{lemma:4.18})
The stabilizer of any point $P$ with segment $(P',P'')$ under the action of 
$\off^{\times}$ via $T_u$ is of the form $U_b$, for some $b\geq 1$. Thus

\begin{equation}
\label{eqn:4.19}
P'=A\sigma T_u(P)=T_uA\sigma (P)=A_u\sigma (P)=P,
%\tag 4.19
\end{equation}
for some $u\in\off$ uniquely determined modulo $U_b$. In other words, if $P$ 
has segment $(P',P'')$ then there is a $u\in \off^{\times}$ such that
$A\sigma T_u (P)=P$ and

\[
\left(
\begin{array}{cc}
u'&0\\0&1 \end{array}
\right),\ \ u'\in U_b,
\]
not only stabilizes $A\sigma (P)$ but also commutes with $T_u, \ u\in
\off^{\times}$. 
\qed

From these facts, Lemma \ref{lemma:4.16} follows. 
As a corollary, we obtain $r(0)=R(0)$,
using Lemma \ref{lemma:3.3} and (\ref{eqn:3.17}).



% === END ORIGINAL adjoint_lift4.tex ===
% === BEGIN ORIGINAL adjoint_lift5.tex ===


\section{The buildings for $PGL(3)$ and $PGL(2)$}
\label{sec:5}

Recall $H_1:=SL(2),\ H_2:=PGL(2),\ G:=PGL(3)$, and $N_1,\ N_2$
denote the Jacquet-Shalika norms (\ref{eqn:1.9}). 
We assume $p\not= 2$ without further mention.
 
Let $h \in N_1(g)\cap H_1(F)$ have eigenvalues $\beta_1,\
\beta_2=\beta_1^{-1}$ and
let $h'\in N_2(g)\cap H_2(F)$ be represented by a matrix $B\in GL(2,F)$
having eigenvalues $\lambda\beta_1,\ \lambda$, for some
$\lambda\in \overline{F}^{\times}$. 
Recall $\bbH1 = \bbh2$ and that the action
of $H_2(F)$ on $\bbh2$ is defined analogously to the action of $G(F)$
on $\bbg$. 
Let $E/F$ denote the extension of $F$ obtained by adjoining the
eigenvalues of $h$; we may assume $E/F$ 
is quadratic since the split case was
handled in \S \ref{sec:2}. 
Since the tree $\bbh2$ has no closed loops, one can easily
show that $h$ has fixed points only if 
$c\beta_i\in \oee^{\times}$, for some $c\in F^{\times}$.
 
Our strategy is
to prove, in the notation of \S \ref{sec:3}, that $s(0)=S(0)$ in the
following manner.

\begin{itemize}
\item[(a)] 
If $|1+\beta_1|_E=1$ then we prove $r(0)=s(0)$, $R(0)=S(0)$.

\item[(b)] 
We prove that the case $E/F$ ramified and $|1+\beta_1|<1$
cannot occur. 

\item[(c)] 
If $|1+\beta_1|_E<1$ then we prove
$s(0)=|\text{Fix}(h')|=1$ and
$S(0)=1$. 

\end{itemize}
Thus, in case (c), we may assume that $E/F$ is
unramified. (That the case (b) cannot occur was
pointed out to me by R. Kottwitz.) 
The $\pm$ sign in (\ref{eqn:3.18}) has been choosen
in such a way that $S(0)\geq 0$ though we won't prove this
until later.

\begin{lemma}
\label{lemma:5.1}
\begin{itemize}
\item[(a)] 
$s(0)\leq r(0)$.

\item[(b)] 
Assume that $|1+\beta_1|=1$. Then the number of
fixed points of $h$ is equal to the number of fixed points of $h'$,
i.e., $r(0)=s(0)$. Moreover, $|S(0)|\leq R(0)=r(0)$.
\end{itemize}
\end{lemma}

\pf
Part (a) follows immediately from Lemma \ref{lemma:3.3}. 
For part (b), recall $\Delta_{H_1}(h):=|1-\beta_1^2||\beta_1|^{-1}
=|1-\beta_1||1+\beta_1||\beta_1|^{-1}$ and
$\Delta_{H_2}(h'):=|1-\beta_1||\beta_1|^{-1/2}$, 
by the definitions in \S \ref{sec:2}.
The lemma now follows immediately from the counting 
formulas (\ref{eqn:3.17})-(\ref{eqn:3.18}) for
$R(0)$ and $S(0)$ and the hypothesis.
\qed

\begin{lemma}
\label{lemma:5.2}
\begin{itemize}
\item[(a)] 
The set {\rm Fix}($A_a\sigma$) is in one--to--one correspondence with the
L's such that

\[
[L] = [\{v\in V\ |\ {^{t}v}\cdot \, ^tA_a^{-1}\cdot v'\in \off ,
\ \forall v'\in L\}],
\]
i.e., the class of $L$ is self--dual with respect to the bilinear form of
$^tA_a^{-1}$.

\item[(b)] 
Assume $|{\rm Fix}(A\sigma )|\not= 0,\
[L]\in {\rm Fix}(A_u\sigma ),\ u\in \off^{\times},$
and $[L]\notin\zz$. Then $u=v^2$, for some $v\in\off^{\times}$, and in
particular $(u,\eps)_2=1$ (in the notation of (\ref{eqn:3.18})).
\end{itemize}
\end{lemma}

\pf
(a) This is a simple calculation using (\ref{eqn:4.3}): 
since $A\sigma$ sends $[L]$
to $[AL^{\vee}]$, $[L]$ is fixed if and only if $L$ belongs to the same
class as

\[
\begin{split}
AL^{\vee}
&=\{Av\in V\ |\ {^{t}v}\cdot v'\in \off ,
\ \forall v'\in L\}\\
&=\{v_o\in V\ |\  v_o=Av{\rm  and }\ {^{t}v}\cdot v'\in \off ,
\ \forall v'\in L\}.
\end{split}
\]
It is easy to see that this is equivalent to the statement in part (a).
\newline
(b) We may rephrase the problem in terms of self--dual lattice
classes. 
Let $P$ be an $^tA^{-1}$--self--dual lattice class. 
By (a) above and Lemma \ref{lemma:4.14},
the transitivity of $T_a$ implies that
there is an $a\in\off^{\times}$ such that $[L]=T_aP$. It follows that
$[L]$ is both
$^tA_u^{-1}$--self--dual and $^tA_{a^{-2}}^{-1}$--self--dual, since $T_a$ and
$A_b\sigma$ commute. 
This implies, by (a), that $[L]$ corresponds to a fixed
point of $T_{ua^2}$. 
However, by Lemma \ref{lemma:4.18} and the discussion following 
(\ref{eqn:4.15}), we know that the stabilizer of 
$T_{ua^2}$ is $U_n$,
where $n\geq 1$ is the distance between $[L]$ and its
projection $[L']$ in $\zz$ (by hypothesis, $[L]\not= [L']$).
Since, by definition,
$U_n \subseteq \{ {\rm squares\ in\ }\off^{\times}\}$, for $n\geq 1$, we
conclude that $u$ is a square in $\off^{\times}$.
\qed

\begin{lemma}
\label{lemma:5.3}
Let $w:=\left(\smallmatrix \ &1\\-1&\ \endsmallmatrix
\right)$, and let $\Lambda_0:=\off\oplus\off\oplus\off$ denote the standard
rank 3 lattice written in terms of a fixed basis for $V_A\oplus W_A$.
The action of $A_a\sigma$ on $\bbg$ is given on $\zz$ by

\[
\begin{split}
\left[ 
\left(
\begin{array}{cc} 
r&\ \\\ &C'
\end{array}\right)
\Lambda_0\right] \longmapsto
&\left[ A_a
\left(
\begin{array}{cc} 
1&\ \\ \ &w
\end{array}\right)
\left(
\begin{array}{cc} 
r^{-1}&\ \\
\ &C'
\end{array}\right)
\left(
\begin{array}{cc} 
1&\ \\\ &w
\end{array}\right)^{-1}\Lambda_0\right]\\
&=\left[ A_a
\left(
\begin{array}{cc} 
1&\ \\\ &w
\end{array}\right)
\left(
\begin{array}{cc} 
r^{-1}&\ \\\ &C'
\end{array}\right)
\Lambda_0\right],
\end{split}
\]
where $r\in F^{\times},\ C'\in {\rm Aut}_F(W_A)\cong GL(2,F)$.
\end{lemma}

\begin{remark} This follows from (\ref{eqn:4.3}) and the fact that
$C^{-1}=(det(C))^{-1}\, ^t(w^{-1}Cw)$. This description of the twisted
action on the building is especially useful when the action of
$A_a\sigma$ is restricted $\zz$.
\end{remark}

\begin{lemma}
\label{lemma:5.4} 
Assume that $N_1(g)$ doesn't contain any
unipotent element. If ${\rm Fix}(A\sigma )$
contains a
vertex \underbar{not} in $\zz$ and if
${\rm Fix}(A\sigma )\cap\zz = \emptyset$ then $|\beta'_1/\beta'_2 +1|<1$,
where $\beta'_i$ are the eigenvalues of $Cw$.
\end{lemma}

\pf 
By hypothesis, we may assume that A is of the form
$\left(\begin{array}{cc} 1 &\ \\ \ &C\end{array}\right)$, 
thanks to Corollary \ref{cor:1.14}.
Suppose that there is an $[L]\in {\rm Fix}(A\sigma )$ not in $\zz$. 
Lemma \ref{lemma:4.16}
implies that we can find $[L'],\ [L'']\in\zz$ forming the characteristic
segment of $[L]$ such that
$[A\sigma L'] = [L''],\ [A\sigma L''] = [L']$. 

Denote by $p',\ p''$
the vertices of $\bbb (SL(2))$ associated to $L',\ L''$,
where $Aw(p')=p'',\ Aw(p'')=p'$.
These are depicted by the lines in Figure \ref{fig:4.8}.
%%%%%%%%%%Last sentence added 3-2003 to avoid redrawing 4.8 
We may consider the midpoint $p_0$ of $p',\ p''$ in the first barycentric
subdivision of $\bbb (SL(2,F))$ as a vertex of the tree over $E$, since
$E/F$ is ramified. By Lemma \ref{lemma:5.3} and a
well--known result on trees, $p_0$ in $\bbb (SL(2,E))$ is a fixed point
of $Cw$ ($Cw$ acts without inversion on the first barycentric subdivision
of $\bbb (SL(2,F))$). By Lemma \ref{lemma:5.3}, this midpoint
is associated to a fixed point $[L_0]$ over E of ($A\sigma$ and)
$A\cdot\left(\begin{array}{cc}
1&\ \\\ &w\end{array}\right)$. By hypothesis 
$[L_0]\notin \zz$.
 
Let
$p_1,\ p_2$ be two neighbors of $p_0$, with $p_i$ a vertex in
$\bbb (SL(2,F))$ and with $Cw(p_1)=p_2,\ Cw(p_2)=p_1$. There is an
apartment containing $p_1,\ p_2$ in which we can write

\begin{equation}
\begin{split}
&p_1 = [\pi^{b_1+1}e_1\off + \pi^{b_2}e_2\off ],\\
&p_2 = [\pi^{b_1}e_1\off + \pi^{b_2}e_2\off ],
\end{split}
\label{eqn:5.5}
\end{equation}
for some $b_i\in\zzz$ and a suitable basis $\{e_1,\ e_2\}$. It follows that
$Cw$ is conjugate over $F$ to
$\left(\smallmatrix u_1\pi^{-1}&u_2\pi^{b_1-b_2}\\ u_3\pi^{b_2-b_1-1}&
u_4\endsmallmatrix\right)$,
for some $u_i\in \off^{\times}$, hence projectively conjugate over $F$ to
$\left(\begin{array}{cc} \ u_1\pi^{b_1-b_2}\ &u_2\pi^{2(b_1-b_2)+1}\\u_3&
u_4\pi^{b_1-b_2+1}\end{array}\right)$. If $b_1-b_2\leq -2$ or
$b_1-b_2\geq 1$ then the quotient of the eigenvalues
of this matrix are equal to $-1 \pm \epsilon$, for some
$|\epsilon |<1$. 
In this case, from Lemma \ref{lemma:4.5} we find that
$Cw$ can have at most one fixed point in the tree over $E$,
and none in the tree over $F$. In this case,
we conclude that
there is no fixed point of $A\sigma$ in $\zz$. The cases $b_1-b_2=-1$
and $b_1-b_2=0$ are similar: if $b_1-b_2=-1$ then $p_1=[\Lambda_0]$
and $p_2=[\pi^{-1}e_1\off +\Lambda_0]$. It follows that
$Cw$ is conjugate over $F$ to
$\left(\smallmatrix u_1&u_2\pi^{-1}\\ u_3&u_4\endsmallmatrix\right)$,
for some $u_i\in \off^{\times}$, and the same conclusion may be drawn.
The case $b_1-b_2=0$ is similar and therefore omitted.  
\qed

\begin{proposition}
\label{prop:5.6}
If $|1+\beta_1|_E = 1$ then $|S(0)| = s(0)$.
\end{proposition}

\pf 
Again, by hypothesis, we may assume that $A$ is of the form
$\left(\begin{array}{cc}
1&\ \\\ &C\end{array}\right)$, thanks to Corollary
\ref{cor:1.14}.
 
Note that by the reduction formulas and Lemma \ref{lemma:5.2}(b), we have

\begin{equation}
R(0) - |S(0)| \leq 2Q(|{\rm Fix}(A\sigma )\cap \zz |
+|{\rm Fix}(A_a\sigma )\cap \zz |),
\label{eqn:5.7}
\end{equation}
where $A$ and $A_a$ represent $\ddd_\sigma (g)$ (see (\ref{eqn:1.25})). 
Thus if ${\rm Fix}(A\sigma )\cap\zz = \emptyset$ were true, 
Lemma \ref{lemma:5.1}
would imply the result. Moreover, since
$|S(0)| \leq R(0) = r(0) = s(0)$, we may immediately dispose of the
case $R(0) = 0$. We may therefore assume that $\Delta (h) < 1$, so that
$|1-\beta_1|<1$ and that $|{\rm Fix}(A\sigma )|>0$.
 
Suppose for the moment that ${\rm Fix}(A\sigma ) \subset \zz$.
By Lemma \ref{lemma:4.16}(b) 
it follows that ${\rm Fix}(A_a\sigma ) =\emptyset $,
so we have $S(0) = R(0)$.
 
Suppose next that ${\rm Fix}(A\sigma )$ is not a subset of $\zz$. Since
${\rm Fix}(A_a\sigma ) \subset \zz$ 
would contradict Lemma \ref{lemma:4.16},
we must then have that ${\rm Fix}(A_a\sigma )$ is not a subset
of $\zz$. 
From Lemma \ref{lemma:4.16} 
we have that ${\rm Fix}(A\sigma ) \cap \zz \not= \emptyset$
or ${\rm Fix}(A_a\sigma ) \cap \zz \not= \emptyset$. Suppose with out loss of 
generality that 
${\rm Fix}(A\sigma ) \cap \zz \not= \emptyset$. 
By (\ref{eqn:1.25}) and 
Lemma \ref{lemma:4.16}(b), if $E/F$ was 
ramified then we may take $a\in \off^{\times}$
(since $F^{\times}/N(E^{\times})$ is represented by a unit in the
case of a ramified extension $E/F$).
In this case, Lemma \ref{lemma:5.2}(b) implies all the signs in 
(\ref{eqn:3.18}) are
the the same, so $|S(0)|=R(0)$. On the other hand, if $E/F$ were
unramified then $a\notin \off^{\times}$ by 
Lemma \ref{lemma:4.14} and then
Lemma \ref{lemma:4.16}(d) forces ${\rm Fix}(A_a\sigma ) \cap \zz = \emptyset$.
Then Lemma \ref{lemma:5.4} implies that $Cw$ has no fixed points, which
by Lemma \ref{lemma:5.3} would contradict our assumption.
\qed

Next we must consider the case where $|1+\beta_1|<1$ and $E/F$ is unramified.
In this case, $\Delta (h') = 1$, so by the formula for the number of fixed
points we have $|{\rm Fix}(h')| = 1$. On the other hand, the same formula
implies that $|{\rm Fix}(h)| > 1$, so it is clear that in this case
$r(0) > s(0)$. We of course want to show that the sum in the
formula for $S(0)$ is equal to 1, so $s(0) = Q$. We want to show that,
if $|1+\beta_1| = q^{-m}$ with $m>0$ and $|{\rm Fix}(A\sigma )|\not= 0$
then

\begin{equation}
S(0):={\pm 2\over{[F^{\times}:U]}}|(1+\beta_1)(1+\beta_2)|^{1/2}
\sum_{a\in F^{\times}/U}
\sum_{\substack{ P\\ {\rm Inv}(A_a\sigma (P),P) = 0}}
\kappa (A_a) = 1.
\label{eqn:5.10}
\end{equation}
In other words, we want to show that

\begin{equation}
{2\over{[F^{\times}:U]}}
\sum_{\substack{ a\in F^{\times}/U\\ (a,\epsilon)_2 =+1}}
\sum_{\substack{ P\\ {\rm Inv}(A_a\sigma (P),P) = 0}} 1
= q^m + {2\over{[F^{\times}:U]}}
\sum_{\substack{ a\in F^{\times}/U\\ (a,\epsilon)_2=-1}}
\sum_{\substack{ P\\ {\rm Inv}(A_a\sigma (P),P) = 0}} 1\ .
\label{eqn:5.11}
\end{equation}
For this it clearly suffices to prove that
\begin{equation}
{2\over{[F^{\times}:U]}}
\sum_{\substack{ a\in F^{\times}/U\\ (a,\epsilon)_2 =+1 }}
\sum_{\substack{ P\\ {\rm Inv}(A_a\sigma (P),P) = 0 }} 1
={{q^{m+1}+q^m-2}\over{q-1}}+R,
\label{eqn:5.12}
\end{equation}
(the ``unramified case''), and

\begin{equation}
{2\over{[F^{\times}:U]}}
\sum_{\substack{ a\in F^{\times}/U\\ (a,\epsilon)_2 =-1 }}
\sum_{\substack{ P\\ {\rm Inv}(A_a\sigma (P),P) = 0 }} 1
={{2(q^m - 1)}\over{q-1}}+R,
\label{eqn:5.13}
\end{equation}
(the ``ramified case'') for some $R$. 
These are due to the fact that every fixed point $P$ of $A$ gives
rise to a characteristic leaf by Lemma \ref{lemma:4.16}, 
hence to a fixed point in the
first barycentric subdivision of ${\cal B}(SL(2,F))$. The difference
of (\ref{eqn:5.12}) and (\ref{eqn:5.13}) 
merely counts the difference between the
sets of fixed points in the 
first barycentric subdivision of ${\cal B}(SL(2,F))$ in the unramified case
and in the ramified case, respectively. 
The number of such fixed 
points is as stated in (\ref{eqn:5.12}) by 
Lemma \ref{lemma:4.5}(a), in the unramified
case, and by Lemma \ref{lemma:4.5}(b) in the ramified case.
 
It remains to dispose of the case $|1+\beta_1|<1$ and $E/F$ ramified.
Although one must be able to deal with this case by a more direct argument,
I will give the following geometric proof. 
Since $|1-\beta_1|=1$ and $E/F$ is ramified, 
it follows from \cite{La3}, p. 50,
for example, that $h'$ has no fixed points. 
This implies that $s(0)=0$ and
(by the proof of Lemma \ref{lemma:5.4}) that
$A\sigma$ has no fixed points in $\zz$. 
Thus each fixed point of $A\sigma$, which we may 
assume exists, is associated to a
characteristic base which is not a vertex.
By the remarks preceeding Figure \ref{fig:4.10}, the fact that
$A\sigma$ has no fixed points in $\zz$ implies
that, in the notation of (\ref{eqn:5.7}),
$A_a\sigma$ will have fixed points in $\zz$ (since $E/F$ is ramified,
there is an element of valuation one in $N_{E/F}(E^\times )$
and therefore we may take $A_a$ to correspond to this element; the midpoint
of the characteristic base isn't a vertex,
whereas the midpoint of the dotted line $D_{A_a}$ associated to the fixed
point $L$ of $A\sigma$ in Figure \ref{fig:4.10} is a vertex in $\zz$
since the valuation of $a$ is one). By Lemma \ref{lemma:5.3}, 
this implies that,
in the notation of the proof of Proposition \ref{prop:5.6},
$Cw$ has fixed points. By Lemma \ref{lemma:5.3},
this contradicts the fact that $A\sigma$ has no fixed points in $\zz$.
 
This completes the proof that $S(0)=s(0)$.

 
% === END ORIGINAL adjoint_lift5.tex ===


\begin{thebibliography}{99}


\bibitem[Bl]{Bl} 
C. Blondel, ``Les repr\'esentations supercuspidales des groupes 
m\'etaplectiques sur GL(2) et leurs caract\'eres'', Bull. Soc. Math. de 
France, (Memoire 18) \underbar{113}(1985)

\bibitem[Bor]{Bor} 
A. Borel, ``Automorphic L--functions'', in Proc. Symp. Pure
Math., vol 33, part 2, AMS, Providence, RI, 1979

\bibitem[C]{C} P. Cartier, ``Representations of p--adic groups'', in Proc. 
Symp. Pure Math., vol 33, part 2, AMS, Providence, RI, 1979

\bibitem[F]{F} Y. Flicker, ``On the symmetric square: orbital integrals'',
Math. Ann. 279(1987)173-191

\bibitem[GJ]{GJ} 
S. Gelbart and H. Jacquet, ``A relation between automorphic forms
on GL(2) and GL(3)'', Ann. Sci. \'Ecole Norm. Sup. \underbar{11}
(1978)471-552

\bibitem[JPSS]{JPSS} H. Jacquet, I. Piatetski--Shapiro, and J. Shalika,
``Automorphic forms on GL(3), I, II'', Ann. Math. \underbar{109} 
(1979) 169--258

\bibitem[KS]{KS} R. Kottwitz and D. Shelstad, in preparation
(since appeared as {\bf Foundations of twisted endoscopy}, 
Asterisque, vol 255, 1999)

\bibitem[Ko1]{Ko1} R. Kottwitz, ``Orbital integrals on GL(3)'', Amer. J. Math.
\underbar{102} (1980) 327--384

\bibitem[Ko2]{Ko2} R. Kottwitz, ``Unstable orbital integrals on SL(3)'', 
Duke Math. J. \underbar{48}(1981)649--664

\bibitem[LL]{LL} J.-P. Labesse and R. P. Langlands,
``L-indistinguishability for
$SL(2)$'', Can. J. Math. \underline{33}(1979)

\bibitem[L]{L} S. Lang, \underbar{Algebra}, $2^{nd}$ ed., Addison-Wesley,
1984

\bibitem[La1]{La1} R. Langlands, ``Some identities for orbital integrals 
attached to GL(3)'', manuscript

\bibitem[La2]{La2} R. Langlands, \underbar{Les Debut d'une Formule des Traces
Stable}, Publ. Math. Univ. Paris VII, 1980 

\bibitem[La3]{La3} R. Langlands, \underbar{Base Change for GL(2)}, 
Ann. Math. Studies, Princeton Univ. Press, 1980 

\bibitem[LS]{LS} R. Langlands and D. Shelstad, ``On the definition of transfer 
factors'', Math. Ann. (1987)219--271

\bibitem[M]{M} I. Macdonald, \underbar{Spherical Functions on a Group of
p--adic Type}, Publ. Ramanujan Inst., no. 2, Madras, 1971

\bibitem[Sh]{Sh} G. Shimura, ``On the holomorphy of certain Dirichlet series'',
Proc. London Math. Soc. \underbar{31} (1975) 79-98

\bibitem[T]{T} J. Tits, ``Reductive groups over local fields'', in Proc. 
Symp. Pure Math., vol 33, part 1, AMS, Providence, RI, 1979

%\bibitem[]{}

%\bibitem[]{} 

%\bibitem[]{} 

\end{thebibliography}

\end{document}






% === SOURCE ARCHIVE: original adjoint_lift_pics.tex ===
% The separate picture study is preserved verbatim after \end{document}.
% These figure drawings already appear, captioned and labeled, in adjoint_lift4.tex.
\documentclass[12pt]{article}
%\usepackage{amssymb}
%\usepackage{amsthm}
%\usepackage{amsgen}
\usepackage{amsfonts}

\newtheorem{theorem}{Theorem}
\newtheorem{corollary}[theorem]{Corollary}
\newtheorem{conjecture}[theorem]{Conjecture}
\newtheorem{lemma}[theorem]{Lemma}
\newtheorem{proposition}[theorem]{Proposition}
\newtheorem{definition}[theorem]{Definition}
\newtheorem{example}[theorem]{Example}
\newtheorem{axiom}{Axiom}
\newtheorem{remark}{Remark}
\newtheorem{exercise}{Exercise}[section]

\def\ppp{{\mathbb{P}}}
\def\hhh{{\mathbb{H}}}
\def\aaa{{\mathbb{A}}}
\def\fff{{\mathbb{F}}}
\def\qqq{\mathbb{Q}}
\def\rrr{\mathbb{R}}
\def\ccc{\mathbb{C}}
\def\zzz{\mathbb{Z}}
\def\nnn{\mathbb{N}}
\def\pf{{\bf proof}:\ }
\def\qed{$\Box$}

\begin{document}


\author{David Joyner}
\title{Adjoint lifting figures}
\date{3-2003}
\maketitle


%%figure (4.10) on page 1121
\begin{picture}(200.00,150.00)(-120.00,0.00)
\thicklines
\put(-50.00,100.00){\line(1,0){250.00}}
\put(-50.00,0.00){\line(1,0){250.00}}
%left leaf
\put(-29.00,50.00){$\bullet$} 
\put(-20.00,-20.00){$L''$} 
\put(-14.00,-2.00){$\bullet$} 
\put(35.00,97.00){$\bullet$} 
\put(-35.00,33.00){$L$} 
\put(40.00,110.00){$L'$} 
\put(-10.00,00.00){\line(3,6){50.00}}

%right leaf
\put(91.00,50.00){$\bullet$} 
\put(100.00,-20.00){$M''$} 
\put(106.00,-2.00){$\bullet$} 
\put(155.00,97.00){$\bullet$} 
\put(86.00,33.00){$M$} 
\put(160.00,110.00){$M'$} 
\put(110.00,00.00){\line(3,6){50.00}}

\thinlines
%right leaf
\put(110.00,00.00){\line(-1,3){17.00}}
\put(92.00,52.00){\line(4,3){65.00}}

%left leaf
\put(-10.00,00.00){\line(-1,3){17.00}}
\put(-28.00,52.00){\line(4,3){65.00}}

%middle line
\put(50.00,00.00){\line(3,6){50.00}}
\put(43.00,33.00){$D_A$} 

\put(-110.00,0.00){$\Lambda'' =\gamma \Lambda''$} 
\put(-110.00,100.00){$\Lambda' =\gamma \Lambda'$} 
\end{picture}

\vskip 1in


%%figure (4.8) on page 1120
\begin{figure}[h!]
\begin{picture}(200.00,150.00)(-120.00,0.00)
\thicklines
\put(-50.00,100.00){\line(1,0){220.00}}
\put(-50.00,0.00){\line(1,0){220.00}}
\put(71.00,50.00){$\bullet$} 
\put(80.00,-20.00){$L''$} 
\put(86.00,-2.00){$\bullet$} 
\put(135.00,97.00){$\bullet$} 
\put(60.00,50.00){$L$} 
\put(140.00,110.00){$L'$} 
\put(90.00,00.00){\line(3,6){50.00}}

\thinlines
\put(90.00,00.00){\line(-1,3){17.00}}
\put(72.00,52.00){\line(4,3){65.00}}

\end{picture}
\vskip .3in
\caption{A characteristic leaf \label{fig:4.8}}
\end{figure}

\vskip 1in

This is figure \ref{fig:4.8}.

\vskip 1in

%%figure (4.6) on page 1119
\begin{picture}(200.00,150.00)(-120.00,0.00)
\thicklines
\put(-80.00,50.00){\line(1,0){220.00}} 
\put(-10.00,60.00){$L_{(-1)}$} 
\put(-2.00,47.00){$\bullet$} 
\put(0.00,50.00){\line(1,-1){50.00}}
\put(100.00,50.00){\line(1,-1){30.00}}
\put(-60.00,-23.00){$L'_{(-1)}$} 
\put(-52.00,-3.00){$\bullet$} 
\put(0.00,50.00){\line(-1,-1){50.00}}
\put(47.00,-3.00){$\bullet$}  
\put(90.00,60.00){$L_{(0)}$} 
\put(-80.00,0.00){\line(1,0){220.00}}
\put(40.00,-23.00){$L'_{(0)}$} 
\put(50.00,0.00){\line(1,1){50.00}}
\put(-50.00,0.00){\line(-1,1){20.00}}
\put(98.00,47.00){$\bullet$}  
\put(158.00,43.00){$\Lambda$}  
\put(158.00,-7.00){$\Lambda'$}  
\end{picture}

\vskip 1in


%%figure (4.17) on page 1123
\begin{picture}(200.00,150.00)(-120.00,0.00)
\thicklines
\put(-100.00,50.00){\line(1,0){270.00}} 
\put(-10.00,60.00){$Q'$} 
\put(-2.00,47.00){$\bullet$} 
\put(0.00,50.00){\line(1,-1){50.00}}
\put(100.00,50.00){\line(1,-1){50.00}}
\put(140.00,-23.00){$Q''$} 
\put(148.00,-3.00){$\bullet$} 
\put(0.00,50.00){\line(-1,-1){30.00}}
\put(47.00,-3.00){$\bullet$}  
\put(90.00,60.00){$P'$} 
\put(-100.00,0.00){\line(1,0){270.00}}
\put(40.00,-23.00){$P''$} 
\put(50.00,0.00){\line(1,1){50.00}}
\put(150.00,0.00){\line(1,1){20.00}}
\put(98.00,47.00){$\bullet$}  
\put(-150.00,20.00){$P',P'',Q',Q''\in {\cal Z}$}  
\end{picture}

\vskip 1in

\begin{thebibliography}{99}


%\bibitem[]{}

%\bibitem[]{} 

%\bibitem[]{} 

\end{thebibliography}

\end{document}
